\chapter{Methodology}
\label{ch:methods}

This chapter defines the methods that the thesis compares.
\S\ref{sec:methodology_problem} fixes the setting: the taxonomy, the score
conventions, the training objective, and the lexicographic problem that the rest
of the chapter addresses.

Three mechanisms are then developed, one per section. All three pursue the same
goal, that the coarse levels of the taxonomy take precedence over the finer ones,
but each enforces it on a different object:

\begin{itemize}
    \item the \textbf{Hierarchical Constraint Cascade} (HCC) constrains the
    \emph{outputs}. The coarse scores are produced first, and the finer ones are
    then obtained as the solution of a linear problem that forces them to agree
    with the coarse ones. (\S\ref{sec:methodology_hcc}).
    \item \textbf{Lexicographic mode} constrains the \emph{update}. Gradients are
    computed as usual, and the contribution of each level is then projected so
    that, to first order, it cannot undo what the levels above it are asking for
    (\S\ref{sec:methodology_lex}).
    \item The \textbf{LH-projection} constrains \emph{what each level learns
    from}. At the point where a level head reads the shared representation, the
    learning signal that head sends back is filtered so that it carries no
    component along the directions in which the heads above it currently read
    (\S\ref{sec:methodology_lhproj}).
\end{itemize}

The first acts in the forward pass and therefore also at inference; the other two
act only in the backward pass, and change nothing about the function a trained
network computes. That asymmetry recurs throughout the chapter and is what makes
the three comparable rather than interchangeable.

Two further sections complete the apparatus.
\S\ref{sec:methodology_inference} fixes how a trained network is turned into a
predicted path, introducing multiple inference methods. \S\ref{sec:methodology_axis} compares
the three mechanisms.

\section{Problem formulation}
\label{sec:methodology_problem}

This section states the problem addressed by the rest of the chapter and fixes
the notation used for it throughout the thesis. The setting is the one already
introduced in \cref{sec:hc_formulation}: in the terms of the Silla and Freitas
classification of hierarchical problems~\cite{silla2011survey}, it is the case
$\langle T, S, FD\rangle$, that is, a tree-shaped hierarchy in which every
instance belongs to exactly one root-to-leaf path and is labelled down to a leaf.
What follows restates that setting in the notation the constructions of this
chapter require, and ends with the optimisation problem they are meant to solve.

\subsection{Taxonomy}
\label{ssec:problem_taxonomy}

A \emph{taxonomy} is a rooted tree $\mathcal{T}$ whose nodes are partitioned into
$L$ levels, ordered from the most general to the most specific and indexed
$1\prec 2\prec\dots\prec L$. Level $l$ contains $C_l$ classes, and
$N:=\sum_{l=1}^{L}C_l$ denotes the total number of nodes. Two structural
assumptions are made throughout:

\begin{enumerate}
    \item \textbf{Unique parent.} Every class at level $l+1$ has exactly one
    parent at level $l$. The parent map is written
    $\pi_l:\{1,\dots,C_{l+1}\}\to\{1,\dots,C_l\}$, and its matrix form is the
    incidence matrix $M_{l,l+1}\in\{0,1\}^{C_l\times C_{l+1}}$ with
    $M_{l,l+1}[i,j]=1$ if and only if $\pi_l(j)=i$. Uniqueness of the parent is
    exactly the statement that every column of $M_{l,l+1}$ contains a single
    non-zero entry.
    \item \textbf{No childless parent.} Every class at level $l$ has at least one
    child at level $l+1$, so that each row of $M_{l,l+1}$ is non-zero. As shown in
    \cref{lemma:diag}, this is precisely the condition under which the HCC
    projection is well defined.
\end{enumerate}

Under these assumptions each leaf determines a unique root-to-leaf path, so a
label may equivalently be given as a leaf identifier or as the path
\[
y=(y_1,\dots,y_L),\qquad y_l\in\{1,\dots,C_l\},\qquad y_l=\pi_l(y_{l+1}).
\]
A path satisfying $y_l=\pi_l(y_{l+1})$ for all $l$ is called \emph{consistent};
the datasets provide consistent ground-truth paths by construction, whereas a
model's prediction need not be consistent, which is what the metrics of
\S\ref{ssec:inference_metrics} quantify. Note that a model could be consistent without 
being correct.

All datasets used in this thesis are given a three-level hierarchy, $L=3$. The
methods are defined for arbitrary $L$ and the restriction is a property of the
experimental setting, not of the formulation; the two exceptions --- the HRN
baseline and the lexicographic mode of \S\ref{sec:methodology_lex}, both of which
require exactly three levels as realised here --- are noted where they
arise.

\subsection{Classifier and score conventions}
\label{ssec:problem_scores}

A hierarchical classifier is a map $x\mapsto\big(z_1,\dots,z_L\big)$
parameterised by $\theta$, with
\[
z_l\;:=\;z_l(x;\theta)\in\mathbb{R}^{C_l}
\]
the vector of scores assigned to the classes of level $l$. The explicit
dependence on $x$ and $\theta$ is omitted whenever it is not needed. Levels are
always indexed coarse-to-fine, so that $z_1$ is the coarsest and $z_L$ the
finest.

Two conventions for producing $z_l$ occur among the model families compared here,
and the distinction matters when the constructions of the following sections are
applied.

\begin{itemize}
    \item \textbf{Per-level heads.} Most families (H-CAST, LH-DNN, HRN,
    HT-CapsNet) emit one score vector per level directly, as the output of a
    level-specific head applied to a shared representation.
    \item \textbf{Taxonomy-node scores.} Hier-COS instead emits a single node
    vector $u\in\mathbb{R}^{N}$ over the whole taxonomy and reduces it to level
    scores by subspace norms,
    \[
    z_l[c]=\big\lVert \Pi_{S_c}\,u\big\rVert_2,
    \]
    where $S_c$ is the taxonomy subspace assigned to node $c$ of level $l$. This
    reduction is non-linear, and consequently the affine constructions of
    \S\ref{sec:methodology_hcc} and \S\ref{sec:methodology_lhproj} act on the
    linear per-level node blocks of $u$ rather than on the subspace norms
    themselves. Where this distinction affects the interpretation of a result it
    is stated explicitly.
\end{itemize}

\subsection{Objective}
\label{ssec:problem_objective}

Unless a section states otherwise, training minimises a scalarised objective of
the form
\[
\mathcal{L}(\theta)
=
\sum_{l=1}^{L} w_l\,\ell_l\big(z_l,y_l\big)
\;+\;
\lambda\,\mathcal{R}\big(z_1,\dots,z_L\big),
\]
where $\ell_l$ is a level-wise loss (cross-entropy, or a family-specific
surrogate such as the capsule margin loss of HT-CapsNet), $w_l>0$ are level
weights with $\sum_l w_l=1$, and $\mathcal{R}$ is an optional family-specific
regulariser with coefficient $\lambda$.

The level weights are a genuine degree of freedom rather than an implementation
detail, and because several claims of this thesis turn on them they are stated
in closed form rather than by name.

\paragraph{Depth-derived weights.} The default weighting of the Hier-COS family,
denoted \texttt{kl\_leaf}, assigns to level $l$ the value $\exp\!\big(1/(L+1-l)\big)$,
normalises the resulting vector in $\ell_2$ and squares it. Since $\ell_2$
normalisation followed by squaring is algebraically identical to $\ell_1$
normalisation of the squares, the resulting weight has the closed form
\[
w_l
=
\frac{\exp\!\big(2/(L+1-l)\big)}{\sum_{k=1}^{L}\exp\!\big(2/(L+1-k)\big)}
=
\operatorname{softmax}_l\!\Big(\tfrac{2}{L+1-l}\Big) ,
\]
so that the squaring acts as an inverse temperature on the weight distribution:
the weights actually used correspond to $\exp\!\big(2/(L+1-l)\big)$ and not, as the
expression $\exp\!\big(1/(L+1-l)\big)$ alone would suggest, to a milder
concentration. This doubling originates in the released implementation of the
method rather than in this work, and is reproduced here deliberately so that the
baseline is the published one; it is recorded because the resulting weight vector
is markedly more leaf-heavy than the written form implies. Two further settings
are obtained from it: \texttt{equal}, with $w_l=1/L$, and \texttt{kl\_coarse},
which reverses the \texttt{kl\_leaf} vector and therefore concentrates weight on
the coarsest level. For $L=3$ the three vectors are, coarse to fine,
$(0.162,\,0.225,\,0.613)$, $(1/3,\,1/3,\,1/3)$ and $(0.613,\,0.225,\,0.162)$.

\paragraph{Branching-derived weights.} The weightings above depend only on the
depth of the taxonomy: they assign the same numbers to a hierarchy that fans out
sharply at the last level and to one that does not. Two alternatives are
therefore defined, in which the weight of a level is a function of the taxonomy
itself rather than of its index. Writing $C_l$ for the number of classes at level
$l$, \emph{cumulative branching} sets
\[
w_l \;\propto\; C_l^{\beta},
\qquad\text{equivalently}\qquad
w = \operatorname{softmax}\big(\beta \log C\big),
\]
with $\beta\ge 0$ (default $\beta=1/2$) controlling how strongly the weighting
follows the class counts: $\beta=0$ recovers \texttt{equal}, and increasing
$\beta$ concentrates weight on the more populated levels. \emph{Marginal
branching} instead takes the coarsest level as the reference unit and scores each
subsequent level by its average expansion with respect to its predecessor,
\[
s_1 = 1,
\qquad
s_l = \frac{C_l}{C_{l-1}}\ \ (l>1),
\qquad
w_l = \frac{s_l}{\sum_k s_k} ,
\]
so that a level receives weight in proportion to the number of decisions it
actually introduces. Both are available for Hier-COS. They
exist because a weighting justified only by level index is difficult to defend as
a modelling choice, whereas both alternatives above are functions of the
taxonomy's own branching structure and can be stated before any result is seen.

\medskip

Because two of the mechanisms introduced below either bypass $w_l$ or replace it
by a uniform weighting, $w_l$ is treated as a \emph{controlled variable} of the
experimental design and is reported for every run; the resulting confound between
arms is discussed in \S\ref{ssec:lex_setting} and in the threats to validity.

\subsection{The lexicographic problem}
\label{ssec:problem_priority}

The objective above is the one the baseline families minimise, and it treats the
$L$ levels as commensurable: a weight $w_l$ prices one level against another, so
a gain at one level can always be bought with a loss at some other. The problem
this thesis addresses is instead the lexicographic one, in which the levels are
ordered by priority and the ordering is absolute:
\begin{equation}
\lexmin_{\theta}\ \ \ell_1(z_1,y_1),\ \ell_2(z_2,y_2),\ \dots,\ \ell_L(z_L,y_L),
\label{eq:hc_lmop}
\end{equation}
that is, the sequence of programs of \cref{def:lmop}, in which level $l+1$ is
minimised only over the minimisers of the levels $1,\dots,l$. The ordering
adopted throughout is $\ell_1\prec\ell_2\prec\dots\prec\ell_L$, coarse before
fine.

The two formulations are not two ways of writing the same thing. As established
in \cref{sec:lexico-real}, no positive real weight vector encodes a lexicographic
order (\cref{prop:debreu}), and for any fixed $w$ there are level-wise changes
under which the scalarised objective improves while the lexicographic order
worsens (\cref{cor:compensation}). The weights of \S\ref{ssec:problem_objective}
therefore express a trade-off, not a priority.

\paragraph{Why coarse before fine.} The ordering in \eqref{eq:hc_lmop} is an
assumption rather than a consequence of the taxonomy, and the argument for it is
the graded cost of errors already discussed in \cref{sec:hc_formulation}. Under a
flat objective, predicting a different breed of the correct bird and predicting a
vehicle are the same event; a hierarchy exists precisely because they are not.

A deep error is therefore preferable to a shallow one, and that is the sense in
which the coarse levels come first. The argument, however, says only which errors
are less damaging when they occur; it does not say that imposing the ordering
during training will produce fewer of the damaging ones. Two qualifications
follow. First, the ordering is a hypothesis about \emph{which} objective should
dominate, and the thesis tests it
rather than assuming it: \S\ref{ssec:lex_modes} applies the
same operator with the priority reversed (fine before coarse), and any effect reproduced under it is
not an effect of the hierarchy (\S\ref{ssec:design_controls}). Second, the
ordering is not free: protecting the coarse level can only restrict the fine
level, so a cost in fine-level accuracy is the expected price rather than a
failure of the method, and the results are read as a position on an
accuracy--consistency trade-off rather than as a uniform improvement.

\section{Hierarchical Constraint Cascade}
\label{sec:methodology_hcc}

The \emph{Hierarchical Constraint Cascade} (HCC) is a differentiable method for enforcing hierarchical relations among the outputs of a multi-level classifier. The method builds on the \textsc{HardNet} framework for imposing hard constraints on neural networks \cite{min2024hardnet}, of which it constitutes a specific instance: the equality-only case, applied recursively along a taxonomy.

This section first presents the general framework (\S\ref{ssec:hardnet}), then formally derives the hierarchical instantiation underlying HCC (\S\ref{ssec:hardnet_to_hcc}), relates the resulting construction to the lexicographic approaches that form the backbone of this thesis (\S\ref{ssec:lexicographic}), and finally discusses the extensions introduced here with respect to the original method (\S\ref{ssec:hcc_extensions}).

\subsection{The HardNet framework}
\label{ssec:hardnet}

\subsubsection{Motivation}

Prior knowledge about the input--output relations of a model is typically imposed in one of two ways. The first, \emph{soft constraints}, penalises violations through a regularisation term in the objective; this offers no guarantee of constraint satisfaction, particularly for inputs far from the training distribution, a condition that is critical in safety-critical applications \cite{min2024hardnet}. The second, \emph{hard constraints} imposed naively --- for instance by redefining the architecture so that the constraint holds by construction through an \emph{ad hoc} parameterisation --- risks a drastic reduction of the model's representational capacity.

HardNet was proposed as a resolution of this trade-off \cite{min2024hardnet}: a framework for constructing networks that satisfy \emph{input-dependent} hard constraints by construction, without sacrificing model capacity, and that remain trainable end-to-end with standard unconstrained optimisation algorithms.

\subsubsection{Definition}

Let $f_\theta:\mathcal{X}\to\mathbb{R}^{n_{\mathrm{out}}}$ be a neural network parameterised by $\theta$. HardNet enforces the input-dependent affine constraint

\[
b^l(x)\le A(x)f(x)\le b^u(x)\qquad\forall x\in\mathcal{X},
\]

with $A(x)\in\mathbb{R}^{n_c\times n_{\mathrm{out}}}$, by appending to $f_\theta$ a differentiable \emph{enforcement layer}, that is, a projection $\mathcal{P}$ such that $\mathcal{P}(f_\theta):\mathcal{X}\to\mathbb{R}^{n_{\mathrm{out}}}$ satisfies the constraint for every $x$ while remaining differentiable with respect to $\theta$ through $f_\theta$. The layer is agnostic to the upstream architecture: it can be attached to any network exposing an output vector of dimension $n_{\mathrm{out}}$ \cite{min2024hardnet}.

Under the assumptions that, for every $x$, the constraint is feasible and $A(x)$ has full row rank, the closed-form solution \textsc{HardNet-Aff} is given by \cite{min2024hardnet}

\[
\mathcal{P}(f_\theta)(x)
=
f_\theta(x)
+
A(x)^{+}
\Big[
\operatorname{ReLU}\!\big(b^l(x)-A(x)f_\theta(x)\big)
-
\operatorname{ReLU}\!\big(A(x)f_\theta(x)-b^u(x)\big)
\Big],
\]

where $A(x)^{+}:=A(x)^{\top}\big(A(x)A(x)^{\top}\big)^{-1}$ denotes the Moore--Penrose pseudoinverse. Two properties established in \cite{min2024hardnet} are directly relevant here:

\begin{itemize}
    \item \textbf{Exact constraint satisfaction}: for each row $a_i$ of $A(x)$, the projected output satisfies $a_i^{\top}\mathcal{P}(f_\theta)(x)=b^l_{(i)}(x)$ or $b^u_{(i)}(x)$ whenever the constraint is violated, and is left unchanged otherwise (Proposition 4.4 in \cite{min2024hardnet}).
    \item \textbf{Preservation of universal approximation}: if the function class realised by $f_\theta$ universally approximates a class $\mathcal{F}$, then the class of projected functions $\mathcal{P}(f_\theta)$ universally approximates the subset of $\mathcal{F}$ satisfying the constraint (Theorem 4.6 in \cite{min2024hardnet}). Enforcing the constraint by construction therefore does not, in the limit, reduce the representational capacity of the model.
\end{itemize}

\subsubsection{Rationale for the choice of HardNet}

Relative to the alternatives surveyed in \cite{min2024hardnet} --- soft constraints (no guarantee), iterative correction procedures such as DC3 (no guarantee within a fixed number of iterations, higher computational cost, sensitivity to the hyperparameters of the correction procedure) and parameterisations of the output set alone, which are input-independent --- HardNet-Aff is presented as the first method with a closed-form (hence inexpensive) forward pass capable of enforcing multiple \emph{input-dependent} affine constraints with a guarantee of exact satisfaction and without a proven loss of expressivity.

This combination of properties --- low computational cost, exact guarantee, architecture independence, and a proven preservation of model capacity --- is the reason this thesis adopts HardNet-Aff as the underlying mechanism for enforcing hierarchical consistency, rather than a soft penalty in the loss or an iterative correction procedure.

\subsection{Notation}
\label{ssec:formulation}

The notation of \S\ref{sec:methodology_problem} is used throughout: $z_l$ denotes
the level-$l$ logits of a classifier parameterised by $\theta$, levels are ordered
$1\prec 2\prec\dots\prec L$ from the most general to the most specific, and
$M_{l,l+1}\in\{0,1\}^{C_l\times C_{l+1}}$ is the incidence matrix of the
parent--child relation between consecutive levels, with exactly one non-zero entry
per column. Only the two structural assumptions of
\S\ref{sec:methodology_problem} --- unique parent, and no parent without children
--- are required by the derivation that follows.

\subsection{From the HardNet-Aff projection to the hierarchical cascade}
\label{ssec:hardnet_to_hcc}

HCC instantiates HardNet-Aff separately for each level transition $l\to l+1$, setting $A(x)\equiv M_{l,l+1}$ --- a \emph{fixed} matrix, not input-dependent, since it encodes the structure of the taxonomy rather than the input --- and imposing an equality-only constraint, that is $b^l=b^u=\hat z_l$:

\[
M_{l,l+1}\hat z_{l+1}=\hat z_l.
\]

An equality constraint is the degenerate case of the box constraint $b^l\le A z\le b^u$ obtained by setting $b^l=b^u$ (Remark 4.1 in \cite{min2024hardnet}). The following paragraph makes explicit the form taken by the general HardNet-Aff projection in this case.

It should be stressed that the constraint is imposed in \emph{logit space}. In general it does not imply the equality
\[
\operatorname{softmax}(\hat z_l)
=
M_{l,l+1}\operatorname{softmax}(\hat z_{l+1}).
\]
HCC must therefore not be read as a normalisation of probability mass between parent and children, but as a structural affine constraint on the representations fed to the objective function.

\subsubsection{Reduction to the equality-only case}

\begin{proposition}
For $b^l=b^u=b$, the HardNet-Aff projection reduces to
\[
\mathcal{P}(f_\theta)(x)=f_\theta(x)-A(x)^{+}\big(A(x)f_\theta(x)-b(x)\big).
\]
\end{proposition}
\begin{proof}
Let $v:=b-Af_\theta(x)$. For every scalar $v$ the identity $\operatorname{ReLU}(v)-\operatorname{ReLU}(-v)=v$ holds, as verified by cases: if $v\ge 0$ then $\operatorname{ReLU}(v)-\operatorname{ReLU}(-v)=v-0=v$; if $v<0$ then $=0-(-v)=v$. The identity extends componentwise to the vector case. Substituting into the general HardNet-Aff expression with $b^l=b^u=b$ yields $\operatorname{ReLU}(b-Af_\theta)-\operatorname{ReLU}(Af_\theta-b)=b-Af_\theta$, from which the claim follows.
\end{proof}

Applying this result with $A=M_{l,l+1}$, $f_\theta(x)=z_{l+1}$ and $b=\hat z_l$ recovers exactly the Euclidean projection onto the affine subspace $\{q:M_{l,l+1}q=\hat z_l\}$:

\[
\hat z_{l+1}
=
z_{l+1}
-
M_{l,l+1}^{\top}
\left(
M_{l,l+1}M_{l,l+1}^{\top}
\right)^{-1}
\left(
M_{l,l+1}z_{l+1}-\hat z_l
\right).
\]

By Proposition 4.4 in \cite{min2024hardnet}, specialised to $b^l=b^u$, this projection satisfies the constraint \emph{exactly} for any $z_{l+1}$, irrespective of whether the constraint was initially violated.

\subsubsection{A structural property specific to taxonomies}

Unlike the general case treated in \cite{min2024hardnet}, in which $A(x)$ may vary with the input and the invertibility of $A(x)A(x)^{\top}$ must be verified pointwise for every $x$, in HCC the matrix $M_{l,l+1}$ is fixed and its "single parent per child" structure admits a closed-form characterisation of $M_{l,l+1}M_{l,l+1}^{\top}$.

\begin{lemma}
\label{lemma:diag}
If every column of $M_{l,l+1}$ contains exactly one non-zero entry, then $M_{l,l+1}M_{l,l+1}^{\top}$ is diagonal, with $i$-th diagonal entry equal to the number of child classes of class $i$ at level $l$.
\end{lemma}
\begin{proof}
Write $M:=M_{l,l+1}$. For $i=i'$ we have $(MM^\top)_{i,i}=\sum_j M[i,j]^2=\sum_j M[i,j]$, since $M$ is binary, which equals the number $n_i$ of children of class $i$. For $i\ne i'$ we have $(MM^\top)_{i,i'}=\sum_j M[i,j]M[i',j]$; by hypothesis, for each column $j$ at most one of $M[i,j]$ and $M[i',j]$ is non-zero, so every summand vanishes and the entry is zero.
\end{proof}

It follows that $M_{l,l+1}M_{l,l+1}^{\top}=\operatorname{diag}(n_1,\dots,n_{C_l})$, which is invertible if and only if every class at level $l$ has at least one child at level $l+1$ --- a structural condition of the taxonomy, verifiable once when the hierarchy graph is built, rather than for every $x$ as required in the general HardNet-Aff setting. Moreover, since $M_{l,l+1}M_{l,l+1}^\top$ is diagonal, the numerically regularised term $M_{l,l+1}M_{l,l+1}^{\top}+\varepsilon I=\operatorname{diag}(n_1+\varepsilon,\dots,n_{C_l}+\varepsilon)$ remains diagonal, so the linear system solved at each forward pass is a diagonal one: it is exact, costs $O(C_l)$ per sample, and could in principle be inverted once before training, since the matrix depends on the taxonomy and not on the input. This is a computational advantage specific to the taxonomic case relative to the general variable-$A(x)$ setting in \cite{min2024hardnet}.

Note further that, for equality-only constraints, feasibility (Assumption 4.3-i in \cite{min2024hardnet}) is implied by full row rank (Assumption 4.3-ii): a surjective linear map admits a solution for any right-hand side. The only condition to be verified on the taxonomy is therefore the absence of parent classes without children.

\subsubsection{Recursive cascade}

Setting $\hat z_1=z_1$ and applying the projection recursively, with the stop-gradient operator $\operatorname{sg}(\cdot)$ on the right-hand side (discussed in \S\ref{ssec:hcc_extensions}), yields

\begin{align}
\hat z_1 &= z_1,\\
\hat z_2 &= P_{1,2}\big(z_2;\operatorname{sg}(\hat z_1)\big),\\
&\;\vdots\\
\hat z_L &= P_{L-1,L}\big(z_L;\operatorname{sg}(\hat z_{L-1})\big),
\end{align}

where each $P_{l,l+1}$ is, by construction, an instance of equality-only HardNet-Aff applied to the level pair $(l,l+1)$. The coarsest level is left unchanged and acts as the reference; each subsequent level is projected onto the affine set determined by the level above it, producing a coarse-to-fine cascade of hierarchical projections.

\subsection{Relation to lexicographic approaches}
\label{ssec:lexicographic}

Since this thesis is concerned with lexicographic formulations of hierarchical classification, the position of HCC within that framework requires explicit justification: HCC is not a lexicographic optimisation method, yet it realises --- in output space and in closed form --- the same priority structure that lexicographic formulations realise in objective space. This subsection makes the correspondence precise, together with its limits.

\subsubsection{Lexicographic priority as level-set restriction}

Given objectives ordered by priority $f_1\prec f_2\prec\dots\prec f_L$, lexicographic optimisation solves the sequence of problems

\[
\min_{q} f_k(q)
\qquad\text{subject to}\qquad
f_j(q)=f_j^{\star},\quad j<k,
\]

where $f_j^\star$ denotes the optimal value attained at stage $j$. The defining feature is not the ordering as such, but the fact that a lower-priority objective may act \emph{only within the level set already fixed by the higher-priority ones}. Equivalently: no finite weight can purchase a trade-off of a higher priority against a lower one, which is precisely what distinguishes lexicographic formulations from the weighted-sum scalarisations used by multi-level architectures such as B-CNN \cite{zhu2017bcnn}. This is the structure exploited by the lexicographic hybrid deep neural network of \cite{fiaschi2024informed}, which combines lexicographic multi-objective optimisation with non-standard analysis \cite{benci2022nonstandard,benci2019measure} to obtain a hierarchical classifier equipped with projection operators before each output layer.

\subsubsection{The structural correspondence}

The HCC projection derived in \S\ref{ssec:hardnet_to_hcc},

\[
\hat z_{l+1}
=
\arg\min_{q\in\mathbb{R}^{C_{l+1}}}
\tfrac{1}{2}\lVert q-z_{l+1}\rVert_2^2
\quad\text{subject to}\quad
M_{l,l+1}q=\hat z_l,
\]

has the same algebraic form. The coarse level contributes an equality constraint --- that is, a level set --- and the fine level is resolved, in the minimum-perturbation sense, strictly inside it. Three consequences follow, and they are the substantive sense in which HCC belongs in a thesis on lexicographic hierarchical classification:

\begin{itemize}
    \item \textbf{One-directional dominance.} The coarse level determines the feasible set for the fine level; the fine level cannot alter the coarse one through the projection branch. The stop-gradient operator extends this asymmetry from the forward pass to the gradient, so that the direction of influence is coarse-to-fine in both.
    \item \textbf{Non-negotiability.} Because the constraint is an equality enforced exactly (Proposition 4.4 in \cite{min2024hardnet}), coarse-level consistency is not traded against fine-level accuracy at any finite value of the level weights $w_l$. In this respect HCC behaves as a lexicographic priority rather than as a scalarised multi-objective compromise, even though the losses themselves are combined additively.
    \item \textbf{Closed form.} The restriction to the higher-priority level set, which in a lexicographic optimisation procedure is obtained by solving a constrained subproblem, is here available in closed form by \cref{lemma:diag}, at the cost of a diagonal rescaling.
\end{itemize}

\subsubsection{Limits of the correspondence}

The analogy must not be overstated. Two differences are substantive and should be borne in mind when interpreting the experimental results.

First, lexicographic optimisation prioritises \emph{objectives} and requires the higher-priority problem to be solved to optimality before the next is considered. HCC constrains \emph{outputs}, and its anchor $\hat z_1=z_1$ is the network's current coarse prediction, not the coarse-optimal one; all levels continue to train jointly. HCC therefore exhibits lexicographic \emph{structure} without lexicographic \emph{sequential optimality}.

Second, the protection afforded by the stop-gradient is partial. It removes the influence of the fine-level loss on the coarse logits along the projection branch, but level-wise losses can still interact through the shared backbone, so parameter-space priority is not enforced.

\subsubsection{Methodological role in this thesis}

Given both the correspondence and its limits, HCC serves three distinct purposes within the present work.

\begin{enumerate}
    \item \textbf{It isolates the contribution of hierarchical priority from that of a particular lexicographic method.} The approach in \cite{fiaschi2024informed} couples lexicographic priority with a specific architecture and with machinery drawn from non-standard analysis. HCC delivers a comparable coarse-to-fine priority structure while remaining architecture-agnostic, allowing the two to be compared and the sources of any observed gain to be attributed.
    \item \textbf{It occupies the hard-constraint end of an axis whose soft end is lexicographic optimisation.} Lexicographic formulations realise priority at optimisation time, with guarantees that hold to first order and per step; HardNet realises it by construction, with exact satisfaction and a proven preservation of expressivity. Comparing the two therefore addresses whether hierarchical priority is better imposed on the optimisation trajectory or on the output contract.
    \item \textbf{It fixes the locus of the intervention while leaving the objective untouched.} HCC adds no term to the loss and no hyperparameter that trades one level against another: it changes only the arguments on which the existing level losses are evaluated. Any difference with respect to the unconstrained baseline is therefore attributable to the constraint itself and not to a reweighting of the objective, which is what makes it comparable, arm for arm, against the two gradient-space mechanisms of \S\ref{sec:methodology_lex} and \S\ref{sec:methodology_lhproj}.
\end{enumerate}

\subsection{Extensions of HCC with respect to HardNet}
\label{ssec:hcc_extensions}

The framework in \cite{min2024hardnet} is designed to impose a constraint --- or a set of constraints --- at a single output point of the network, typically specified by a known function of the problem (for instance a control barrier function, or the constraints of an optimisation problem). HCC adapts it to a different setting: a chain of $L-1$ projections applied in sequence, in which the right-hand side of each projection is not an externally known function but the \emph{projected output of the previous step of the same network}. This difference motivates a design choice that is not part of the original HardNet-Aff formulation, and which should not be attributed to its authors.

\paragraph{Stop-gradient.} In the original formulation the right-hand side $b(x)$ does not depend on $\theta$, so the question of gradient coupling between $b$ and the projection does not arise. In HCC, by contrast, $\hat z_l$ is itself the result of a projection that is differentiable with respect to $\theta$. Without further provision, the level-$(l+1)$ loss would backpropagate through the projection branch and directly modify $\hat z_l$, mixing the objectives of the two levels along that path. The operator $\operatorname{sg}(\cdot)$ --- which returns its argument unchanged in the forward pass and annihilates its gradient in the backward pass --- instead treats $\hat z_l$ as an anchor, realising the coarse-to-fine dependence discussed in \S\ref{ssec:lexicographic}. As noted there, this does not amount to full lexicographic protection of the shared parameters.

The stabilised form of the projection used throughout is therefore

\[
\hat z_{l+1}
=
z_{l+1}
-
M_{l,l+1}^{\top}
\left(
M_{l,l+1}M_{l,l+1}^{\top}
+
\varepsilon I
\right)^{-1}
\left(
M_{l,l+1}z_{l+1}-\operatorname{sg}(\hat z_l)
\right).
\]

For $\varepsilon=0$ and full row rank the affine constraints are satisfied exactly; for $\varepsilon>0$ they are satisfied up to the numerical precision determined by the stabilisation term.

\paragraph{Activation.} The projection is in force from the first optimisation
step: every run reported in this thesis uses $\hat z_l$ in place of $z_l$ for the
whole of training and at evaluation, with no warm start and no gradual
introduction of the constraint. The implementation exposes a blending
coefficient $\alpha\in[0,1]$ that would interpolate between the raw and the
projected logits, and it is logged at every epoch as
\texttt{proj\_constraint\_alpha}; in every configuration used here it is
identically $1$, so the effective logits are $\hat z_l$ throughout and the
coefficient plays no experimental role. It is mentioned only because the
diagnostic is the quantity from which the experimental chapters verify that the
constraint was in fact applied, rather than inferring it from the name of a run.

\paragraph{What the constraint costs the gradient.} Enforcing the constraint from
the first step is not free, and the price is best stated exactly, since
\cite{min2024hardnet} raises it as a hazard --- its Appendix~C reports that an
enforcement layer can produce a vanishing gradient when the loss gradient with
respect to the projected output lies in the span of the active constraint vectors,
a condition that becomes more likely as the number of constraints approaches the
output dimension. Under equality constraints every constraint is active at every
step, so the condition is not occasional but structural, and in the taxonomic case
it takes a closed form.

\begin{proposition}
\label{prop:hcc_backward}
Let $\varepsilon=0$ and write $Q:=M_{l,l+1}^{\top}(M_{l,l+1}M_{l,l+1}^{\top})^{-1}
M_{l,l+1}$. Then $\partial\hat z_{l+1}/\partial z_{l+1}=I-Q$, where $Q$ is the
orthogonal projector onto $\operatorname{row}(M_{l,l+1})$, of rank $C_l$.
Consequently the gradient reaching the raw level-$(l+1)$ logits is confined to a
subspace of dimension $C_{l+1}-C_l$, and $Q$ replaces each coordinate by the mean
of its sibling group, so $I-Q$ centres each sibling group to zero mean.
\end{proposition}
\begin{proof}
The anchor is detached, so differentiating the stabilised form at $\varepsilon=0$
gives $I-M^{\top}(MM^{\top})^{-1}M=I-Q$. $Q$ is symmetric and idempotent with range
$\operatorname{row}(M)$, hence the orthogonal projector onto it, and its rank is the
rank $C_l$ of $M$. By \cref{lemma:diag}, $(MM^{\top})^{-1}=\operatorname{diag}(1/n_i)$,
and row $i$ of $M$ is the indicator of the children of parent $i$; therefore
$(Qv)_j$ is the average of $v$ over the siblings of $j$.
\end{proof}

The fine head therefore never receives the component of its own learning signal
that would raise or lower a whole sibling group together: that degree of freedom
has been delegated to the level above, which is precisely the priority the
mechanism exists to impose. It is the backward-pass counterpart of
\cref{cor:hcc_sibling_shift}, which shows the forward correction to be a common
additive shift per sibling group. The two statements are the same fact seen from
either side of the projection.

Two consequences follow for the interpretation of the results. First, the fraction
of directions removed from level $l+1$ is $C_l/C_{l+1}$, the reciprocal of the
average fan-out at that transition, so the constraint bites hardest wherever the
taxonomy branches least --- and which transition that is depends on the dataset,
not on the depth. It is the coarse-to-middle transition on CIFAR-100 ($8/20$
against $20/100$) and on CUB-200 ($13/38$ against $38/200$), but the
middle-to-fine one on FGVC-Aircraft ($70/100$ against $30/70$), whose taxonomy is
nearly flat at its last level. Any level-wise reading of an HCC result should be
made against the fan-out of the corresponding transition rather than against its
depth. Second, because the constraint is on from the first step, this is a
permanent feature of a run rather than a transient of early training, and no
schedule attenuates it. Whether it is damaging is an empirical question that the
algebra does not settle: the removed component is exactly the one the constraint
declares the finer level has no authority over.

\paragraph{Multi-level cascade.} Finally, whereas HardNet-Aff is applied and evaluated as a single terminal layer in \cite{min2024hardnet}, HCC applies it recursively across $L-1$ successive transitions. The guarantees of exact satisfaction (Proposition 4.4) and preservation of universal approximation (Theorem 4.6) established in \cite{min2024hardnet} concern a single application of the projection; their formal extension by induction over levels to the full cascade with $L>2$ is not proved in this work and remains an open theoretical question. Experimentally this means that exact satisfaction between two consecutive levels is guaranteed by each projection taken individually (\cref{lemma:diag} and the preceding proposition), whereas the effect of composing the projections with the stop-gradient on the overall representational capacity must be assessed empirically.

\subsection{Objective function}

HCC does not introduce a new loss: it modifies the logits on which the objectives already defined by the classifier are evaluated. Writing $\ell_l$ for the loss associated with level $l$, the general objective is

\[
\mathcal{L}_{\mathrm{HCC}}
=
\sum_{l=1}^{L}
w_l\,\ell_l\big(\hat z_l,y_l\big)
+
\lambda\,\mathcal{R}\big(\hat z_1,\dots,\hat z_L\big),
\]

where $w_l$ are the level weights of \S\ref{ssec:problem_objective} and $\mathcal{R}$ an optional model-specific regularisation term. Because the projection is differentiable with respect to the logits, the model does not merely have its outputs corrected after the fact: it is trained on the corrected representations from the first step, and adapts its raw logits to the correction that will be applied to them.

\subsection{Architecture independence}

Like the HardNet enforcement layer, which is presented as applicable to any neural network provided its output contract is specified \cite{min2024hardnet}, HCC is defined on the output contract of the classifier rather than on its internal structure. Only the following are required:

\begin{enumerate}
    \item a vector of differentiable logits for each level of the hierarchy;
    \item a complete taxonomy, with a unique parent for every non-root class and no parent class without children (the condition required by \cref{lemma:diag});
    \item a consistent ordering of the levels from coarsest to finest.
\end{enumerate}

The same cascade can therefore be attached to models with transformer, convolutional, capsule-based or taxonomic-frame backbones, provided their outputs can be reduced to this common interface. The formulation supports an arbitrary number of levels by recursive repetition of the same operation; a given instantiation may restrict itself to a three-level hierarchy without that restriction belonging to the general definition of the method.

This independence concerns the applicability of the mechanism and does not imply that HCC produces the same effect on all architectures. The effectiveness of the constraint, and the interaction between the cascade and the stop-gradient, must be verified experimentally for each combination of model, dataset and decoding strategy.


\section{Gradient-space lexicographic optimisation}
\label{sec:methodology_lex}

Section~\ref{sec:methodology_hcc} imposed hierarchical priority on the \emph{outputs} of the
classifier: the coarse logits define an affine set inside which the fine logits are resolved.
That construction has lexicographic structure, but, as noted in \S\ref{ssec:lexicographic}, it
leaves the shared parameters unprotected --- the level-wise losses still interact freely through
the backbone. This section describes the complementary mechanism developed in this work, which
this thesis refers to as \emph{lexicographic mode}: priority is imposed directly on the
\emph{update direction} in parameter space, by projecting the gradients of lower-priority
objectives so that, to first order, they cannot act against the higher-priority ones.

\subsection{Setting and requirements}
\label{ssec:lex_setting}

Lexicographic mode assumes a classifier exposing exactly three differentiable scalar level
objectives
\[
\ell_1(\theta),\quad \ell_2(\theta),\quad \ell_3(\theta),
\]
associated with the coarse, middle and fine levels and ordered by decreasing priority,
$\ell_1 \prec \ell_2 \prec \ell_3$. Writing
\[
g_l := \nabla_\theta \ell_l(\theta) \in \mathbb{R}^{P}
\]
for the gradient of the $l$-th objective with respect to the vector $\theta$ of trainable
parameters, each optimisation step computes the three gradients separately by three backward
passes through the shared graph, rather than a single backward pass on a scalarised total loss.

Two consequences of this design must be stated explicitly, because they affect the
interpretation of every experiment run in this mode.

\begin{itemize}
    \item \textbf{The scalarised loss is not used for the update.} In lexicographic mode the
    quantity actually applied to the parameters is a combination of the three per-level
    gradients; the weighted total loss is still computed and reported, but never backpropagated.
    Any term of the objective that is not decomposable into the three level losses therefore
    contributes nothing to the update. This is why the mode requires H-CAST to be run without
    its global KL term: that term is not part of any level objective and would be silently
    discarded. The analogous requirements for the baselines are that HRN uses its decomposed
    level-marginal objective, with the leaf cross-entropy folded into the fine level; that
    Hier-COS replaces its published KL objective, which yields a single coupled scalar, by one
    of the two objectives that expose one differentiable term per level; and that HT-CapsNet
    uses its three capsule margin losses. LH-DNN is excluded, since it does not expose three
    independent differentiable level objectives in the required form.
    \item \textbf{The level weights survive in some families and not in others.} The projection
    acts on the three tensors a family publishes as its level objectives, and those tensors do
    not carry the level weights $w_l$ of \S\ref{ssec:problem_objective} uniformly. In H-CAST and
    HT-CapsNet the published level losses are the raw, unweighted ones --- the weights are
    applied only when the scalarised total is assembled, which lexicographic mode never
    backpropagates --- so activating the mode silently replaces the chosen weighting by a
    uniform one. In Hier-COS the level weight is already folded into the cross-entropy part of
    each level loss (though not into the node-norm regulariser), so it survives the projection
    and continues to act. HRN's level-marginal objective exposes no configurable level weights
    at all.

    This is a confound between the lexicographic and non-lexicographic arms wherever the
    baseline uses non-uniform weights, and it is treated as such in the threats to validity: for
    H-CAST and HT-CapsNet the comparison isolates ``priority plus uniform weighting'' against
    ``weighted scalarisation'', not priority alone. It also means that the weighting is not a
    variable one may assume away when reading a lexicographic result, and every run therefore
    reports the weighting it was launched with.

    One consequence of the projection's algebra limits how far the weights can matter. The
    coefficient of \S\ref{ssec:lex_projection} normalises by $\lVert r\rVert_B^2$, so
    $\Pi_B(g;\mu r)=\Pi_B(g;r)$ for every $\mu>0$: rescaling a \emph{reference} gradient leaves
    the projection unchanged. Rescaling a \emph{target} gradient rescales its projection by the
    same factor. The weights therefore act on the magnitudes summed into the final direction,
    and --- through the resultant $g_1+m$ of \S\ref{ssec:lex_modes}, which is not invariant under
    separate rescaling of its two summands --- on the reference direction used for the fine
    level, but never on the pairwise projections taken one at a time.
\end{itemize}

\subsection{Parameter blocks induced by gradient support}
\label{ssec:lex_blocks}

Projecting one gradient onto another is meaningful only on the coordinates where both objectives
actually act. In a hierarchical network the three losses do not reach the same parameters: a
level-specific head is reached by one objective only, while the backbone is reached by all three.
The parameters are therefore partitioned by their \emph{gradient support}
$S(i) := \{\,l : (g_l)_i \not\equiv 0\,\}$, a coordinate counting as outside the support of
$\ell_l$ when the objective $\ell_l$ does not reach it:

\[
\begin{aligned}
T_1 &:= \{\, i : S(i) = \{1,2,3\} \,\} && \text{parameters reached by all three levels,}\\
T_2 &:= \{\, i : S(i) = \{1,2\} \,\}   && \text{reached by coarse and middle only,}\\
T_3 &:= \{\, i : S(i) = \{1\} \,\}     && \text{reached by the coarse level only,}\\
R   &:= \{\, i : \text{otherwise} \,\} && \text{remaining coordinates, typically level-specific heads.}
\end{aligned}
\]

For a block $B$ define the restricted inner product
$\langle u,v\rangle_B := \sum_{i\in B} u_i v_i$ and the induced norm
$\lVert u\rVert_B$. All projection coefficients below are computed with respect to these
restricted inner products, one block at a time, and coordinates in $R$ are left untouched.

The three named blocks are exactly the nested supports
$\{1\}\subset\{1,2\}\subset\{1,2,3\}$, which is the pattern produced by a shared trunk from which
the level heads branch off in coarse-to-fine order --- the structure every family considered here
has. $R$ therefore collects two kinds of coordinate. The first are the level-specific heads,
where the objectives do not share a parameter and there is no conflict to resolve; rescaling a
head gradient by a coefficient estimated elsewhere in the network would be an artefact rather
than a correction. The second are any support patterns outside the nesting, such as $\{2,3\}$,
which would arise only in an architecture where the coarse head branches off before a subtrunk
shared by the two finer levels; such coordinates carry a genuine conflict that this partition
does not address, and their absence in the architectures used here is a property of those
architectures rather than of the construction.

\subsection{The projection operator and the two removal rules}
\label{ssec:lex_projection}

Given a target gradient $g$, a reference gradient $r$ and a block $B$, the basic operation is the
Euclidean projection of $g$ onto the orthogonal complement of $r$, restricted to $B$:

\[
\Pi_B(g;r)_i =
\begin{cases}
g_i - \gamma\, r_i, & i \in B,\\[1ex]
g_i, & i \notin B,
\end{cases}
\qquad
\gamma \;=\; \frac{\langle g,r\rangle_B}{\langle r,r\rangle_B} .
\]

Whether the coefficient $\gamma$ is actually applied is governed by a \emph{removal rule}, and the
choice between the two rules defined below is a genuine methodological decision rather than an
implementation detail: they encode two different readings of what it means for one objective to
be subordinate to another.

\paragraph{\texttt{orthogonalize\_all} (default): subordination.} The component along $r$ is
removed unconditionally, irrespective of the sign of $\langle g,r\rangle_B$, so that
\[
\big\langle \Pi_B(g;r),\, r\big\rangle_B = 0 .
\]
The lower-priority objective is confined to the orthogonal complement of the higher-priority
direction, which is the parameter-space analogue of the level-set restriction of
\S\ref{ssec:lexicographic}: the subordinate objective may act only where the dominant one is, to
first order, indifferent. Removing the component also when the two gradients agree is what
distinguishes subordination from mere de-conflicting: an objective allowed to keep its positive
component along a higher-priority direction is still, in effect, being summed with it, and the
size of its contribution to the dominant objective's progress is then again a matter of relative
magnitude --- exactly the trade-off a lexicographic ordering is meant to forbid.

\paragraph{\texttt{conflict\_only}: de-conflicting.} The coefficient is applied only when the two
gradients disagree, $\langle g,r\rangle_B<-\varepsilon$, and the target is passed through
unchanged otherwise. Writing $\tilde g$ for the result,
\[
\big\langle \tilde g, r\big\rangle_B
=
\begin{cases}
0, & \langle g,r\rangle_B < -\varepsilon,\\[0.6ex]
\langle g,r\rangle_B \;(\ge -\varepsilon), & \text{otherwise},
\end{cases}
\]
so the guarantee becomes one-sided: the lower-priority objective can no longer \emph{oppose} the
higher-priority direction, but it is free to \emph{reinforce} it. This is the gradient-surgery
rule familiar from multi-task learning, where the goal is to remove destructive interference
rather than to install an ordering. It is included here for exactly that contrast. Under
\texttt{orthogonalize\_all} the higher level's first-order progress is a function of its own
gradient alone; under \texttt{conflict\_only} it can still be accelerated by the levels below it,
so the two rules differ not in how much interference they remove but in whether the ordering they
impose is an ordering at all. The comparison between them separates ``priority'' from
``compatibility'' as explanations of any observed effect.

\paragraph{Degeneracy guard.} One guard completes the definition, under both rules. If
$\lVert r\rVert_B^2 \le \varepsilon$ the coefficient is ill-conditioned and the projection is
skipped, $\Pi_B(g;r):=g$, with $\varepsilon=10^{-12}$ throughout. How often the guard fires, and
how often \texttt{conflict\_only} declined to act because the gradients agreed, is recorded at
every step, so that a run can be audited for how often the nominal operator was actually
applied.

\subsection{Composition modes}
\label{ssec:lex_modes}

With three objectives the pairwise operator has to be composed, and the composition fixes what
each level is made subordinate to. Two schemes are defined; in both, the final update direction
passed to the optimiser is the sum $d := \tilde g_1 + \tilde g_2 + \tilde g_3$ of the (possibly
projected) level gradients.

\paragraph{\texttt{coarse\_first} (default).} The coarse gradient is left unchanged; the middle
gradient is orthogonalised against it, separately on $T_2$ and on $T_1$; the fine gradient is
orthogonalised, on $T_1$, against the \emph{resultant} of the two higher-priority directions:
\[
\tilde g_1 := g_1,\qquad
\tilde g_2 := \Pi_{T_2}\!\big(g_2;g_1\big)\ \text{on }T_2,\quad
m := \Pi_{T_1}\!\big(g_2;g_1\big)\ \text{on }T_1,\qquad
\tilde g_3 := \Pi_{T_1}\!\big(g_3;\,g_1+m\big).
\]
The reference for the fine level is $g_1+m$ rather than $g_1$ and $m$ separately, and this is a
substantive choice: what the fine level is forbidden to oppose is the update the two levels above
it have jointly agreed on, not each of their directions in isolation.
\Cref{prop:lex_firstorder} states exactly what that buys and what it does not.

\paragraph{\texttt{fine\_first}.} The priority order is inverted, with the roles of $g_1$ and
$g_3$ exchanged and the same composition applied: the fine gradient is left unchanged, the middle
gradient is orthogonalised against it on $T_1$ (and left unchanged on $T_2$, where the fine level
does not reach), and the coarse gradient is orthogonalised against the raw middle gradient on
$T_2$ and against the resultant $g_3+\tilde g_2$ on $T_1$. This mode is not a hierarchical
hypothesis but a \emph{control}: it applies exactly the same amount of gradient surgery, with the
same operator and the same removal rule, in the opposite order, and so separates the effect of
hierarchical priority from the effect of orthogonalisation as such. Any gain observed under
\texttt{coarse\_first} that is reproduced under \texttt{fine\_first} cannot be attributed to the
hierarchy.

The two modes apply the same composition to the reversed priority order, block by block, over
whichever levels reach the block in question, and each performs the same number of projections:
two on $T_1$ and one on $T_2$. Every statement proved below for \texttt{coarse\_first} therefore
holds for \texttt{fine\_first} after reversing the order --- on $T_1$ with the roles of $g_1$ and
$g_3$ exchanged, and on $T_2$, where the fine level is absent by definition, with the middle level
dominating the coarse one. The blocks themselves are not symmetric, since $T_2$ and $T_3$ are
defined by supports that always contain the coarse level; this asymmetry is a property of where
the level objectives reach, identical in both modes, and not something the reversal introduces.
Only \texttt{coarse\_first} is discussed explicitly from here on.

\subsection{First-order priority guarantee}
\label{ssec:lex_guarantee}

The sense in which this construction realises lexicographic priority is a first-order, per-step
one, and it is not uniform across the parameter blocks. It is exact on $T_2$ and $T_3$, where at
most two levels compete, and weaker on $T_1$, the shared trunk, where all three do. Stating the
two cases separately is not pedantry: $T_1$ is by far the largest block, so it is the weaker
statement that governs how the results should be read.

\begin{proposition}
\label{prop:lex_firstorder}
Assume \texttt{coarse\_first} with the \texttt{orthogonalize\_all} rule and all degeneracy guards
inactive, and let $d$ be the assigned direction.

\emph{(a) On $T_2$}, where $d=g_1+\Pi_{T_2}(g_2;g_1)$ and $m_2:=\Pi_{T_2}(g_2;g_1)$,
\[
\langle d, g_1\rangle_{T_2}=\lVert g_1\rVert^2_{T_2},
\qquad
\langle d, g_2\rangle_{T_2}=\langle g_1,g_2\rangle_{T_2}+\lVert m_2\rVert^2_{T_2}.
\]

\emph{(b) On $T_1$}, where $d=g_1+m+\tilde g_3$,
\[
\text{(i)}\quad \langle d,\, g_1+m\rangle_{T_1}=\lVert g_1\rVert^2_{T_1}+\lVert m\rVert^2_{T_1},
\qquad
\text{(ii)}\quad \langle d, g_1\rangle_{T_1}=\lVert g_1\rVert^2_{T_1}+\rho,
\]
with $\rho:=\langle \tilde g_3, g_1\rangle_{T_1}=-\langle \tilde g_3, m\rangle_{T_1}$, a residual
of arbitrary sign to which the middle level contributes nothing.
\end{proposition}
\begin{proof}
(a) $\langle m_2,g_1\rangle_{T_2}=0$ by construction, which gives the first identity. For the
second, write $g_2=m_2+c\,g_1$ with $c=\langle g_2,g_1\rangle_{T_2}/\lVert g_1\rVert^2_{T_2}$;
then $\langle m_2,g_2\rangle_{T_2}=\lVert m_2\rVert^2_{T_2}$, and adding
$\langle g_1,g_2\rangle_{T_2}$ gives the claim.

(b) By construction $\langle m,g_1\rangle_{T_1}=0$ and, since $\tilde g_3=\Pi_{T_1}(g_3;g_1+m)$,
$\langle \tilde g_3,g_1+m\rangle_{T_1}=0$. For (i), expand
$\langle d,g_1+m\rangle_{T_1}=\lVert g_1+m\rVert^2_{T_1}+\langle \tilde g_3,g_1+m\rangle_{T_1}
=\lVert g_1\rVert^2_{T_1}+\lVert m\rVert^2_{T_1}$, using orthogonality of $g_1$ and $m$. For
(ii), $\langle d,g_1\rangle_{T_1}=\lVert g_1\rVert^2_{T_1}+\langle m,g_1\rangle_{T_1}
+\langle \tilde g_3,g_1\rangle_{T_1}=\lVert g_1\rVert^2_{T_1}+\rho$, and
$\rho=-\langle \tilde g_3,m\rangle_{T_1}$ follows from
$\langle \tilde g_3,g_1\rangle_{T_1}+\langle \tilde g_3,m\rangle_{T_1}=0$.
\end{proof}

Since the directional derivative of $\ell_l$ along the step $\theta\mapsto\theta-\eta d$ is
$-\eta\langle g_l,d\rangle+O(\eta^2)$, the proposition reads as follows.

\begin{itemize}
    \item \textbf{Where only two levels compete, the priority is exact.} By (a), on $T_2$ the
    first-order decrease of the coarse objective is exactly the one its own gradient would
    produce: the middle objective cannot slow it, whatever its magnitude. No finite rescaling of
    the lower level can purchase a trade-off against the higher one --- the defining property of
    a lexicographic, as opposed to a scalarised, formulation. The middle objective in turn
    receives a guaranteed non-negative contribution $\lVert m_2\rVert^2_{T_2}$ from its own
    projected gradient plus a term $\langle g_1,g_2\rangle_{T_2}$ of arbitrary sign contributed by
    the coarse level. Degradation of a level by a \emph{higher}-priority one is permitted;
    degradation by a \emph{lower}-priority one is excluded. The same holds on $T_3$ trivially,
    where the direction is $g_1$ unmodified, and on $R$, where each coordinate keeps its own
    level gradient.
    \item \textbf{On the shared trunk, what is protected is the resultant, not the coarse level
    alone.} By (b)(i), the fine level cannot reduce the first-order progress of the assembled
    higher-priority direction $g_1+m$ below $\lVert g_1\rVert^2+\lVert m\rVert^2$, which is
    strictly positive whenever either gradient is non-degenerate. But by (b)(ii) the coarse level
    taken by itself is not fully protected on $T_1$: the fine level retains a residual $\rho$
    along $g_1$. What (b) does establish is that this residual is \emph{conservative} between the
    two higher levels --- whatever the fine level takes from the coarse objective it returns to
    the projected middle direction, and conversely --- so it can shift first-order progress
    between levels $1$ and $2$ but cannot remove it from the pair. The middle level, by contrast,
    contributes exactly zero to $\langle d,g_1\rangle_{T_1}$: the two-level ordering
    $\ell_1\prec\ell_2$ is exact everywhere, and it is only level $3$ that leaks, and only into
    the $g_1$--$m$ plane.
    \item \textbf{The residual is measured, not assumed away.} $\rho$ is exactly what the
    post-projection cosine between $\tilde g_3$ and $g_1$ on $T_1$ reports. It is expected to be
    non-zero, and it is logged at every epoch (\S\ref{ssec:lex_implementation}); any claim in the
    experimental chapters that the ordering was tight on the trunk is made against that
    measurement rather than against the construction.
\end{itemize}

Under the \texttt{conflict\_only} rule every equality above becomes an inequality in the
favourable direction, up to the guard band: $\langle d,g_1\rangle_{T_2}\ge\lVert
g_1\rVert^2_{T_2}-\varepsilon$, and correspondingly for the resultant on $T_1$. The reason is
that a projection which does not fire leaves an inner product that was, by the rule's own test,
no smaller than $-\varepsilon$. The lower levels can then only help the higher ones, never hinder
them --- a weaker and qualitatively different statement, since it removes interference without
installing an ordering: the size of the help still depends on relative magnitudes, which is
precisely the trade-off a lexicographic ordering is meant to exclude.

Four limitations bound the scope of the guarantee and should be carried into the interpretation
of the results.

\begin{enumerate}
    \item \textbf{First-order and per-step.} The statement concerns directional derivatives at the
    current $\theta$, not the value of the objectives after the step: it guarantees a descent
    \emph{direction} for the dominant objective, not monotone decrease, and says nothing about the
    limit of the trajectory. It is therefore not sequential lexicographic optimality: the coarse
    objective is never solved to optimality before the others are considered. As with HCC, this is
    lexicographic \emph{structure} without lexicographic \emph{sequential optimality}.
    \item \textbf{The optimiser is not gradient descent.} The projected direction $d$ is handed to
    the family's own optimiser in place of the ordinary gradient --- AdamW for H-CAST, Adam for
    HT-CapsNet, SGD with momentum $0.9$ for Hier-COS and HRN --- and none of these applies
    $-\eta d$. Momentum mixes the current direction with past ones, so the applied step is not
    orthogonal to $g_1$ even where $d$ is; adaptive per-coordinate preconditioning additionally
    fails to preserve inner products. \Cref{prop:lex_firstorder} therefore constrains the assigned
    gradient rather than the realised step, and weight decay acts on top of both.
    \item \textbf{Stochasticity.} All gradients are mini-batch estimates. Orthogonality is enforced
    between the estimates, not between the population gradients, and the projection coefficients
    inherit their variance.
    \item \textbf{Block restriction.} The guarantee holds inside each block. It says nothing about
    coordinates in $R$, where the level objectives do not compete, nor about interactions mediated
    by the forward pass rather than by the shared parameters.
\end{enumerate}

\subsection{Cost and observable quantities}
\label{ssec:lex_implementation}

\paragraph{Cost.} Each optimisation step computes three backward passes over the shared graph
instead of one, and holds three gradients of the trainable parameters at once. The same three
per-level gradients are what the diagnostics below are computed from, so they are obtained once
and reused.

\paragraph{What is observed.} The guarantee of \S\ref{ssec:lex_guarantee} is conditional: it
holds under the \texttt{orthogonalize\_all} rule, only where the degeneracy guard did not fire,
and --- on the shared trunk --- only for the resultant rather than for the coarse level alone.
Whether its hypotheses were met, and how large the residual $\rho$ was, is therefore a matter of
measurement rather than of intention, and four families of quantity are recorded at every epoch
for that purpose. The pre-projection cosines between level gradients, on each block, measure how
much conflict there is to resolve. The post-projection cosines verify which orthogonalities were
in fact achieved --- the one against the resultant is expected to vanish, the one against $g_1$
alone is not, and the gap between them is $\rho$. The per-block gradient norms before and after
projection quantify how much of the lower-priority signal was removed. Records of whether each
projection was applied or skipped establish how often the nominal operator was the one actually
used, which matters especially under \texttt{conflict\_only}, where declining to act is the
common case. Claims about lexicographic behaviour in the experimental chapters rest on these
measurements rather than on the setting a run was launched with.

\subsection{Position relative to HCC}
\label{ssec:lex_vs_hcc}

Lexicographic mode and HCC realise the same coarse-to-fine ordering at two different loci.
HCC constrains the logits: the restriction is exact,
holds per sample, and holds at inference as well as during training, but it governs only what
the heads emit and leaves the shared parameters free to be pulled in incompatible directions by
the level losses. Lexicographic mode constrains the assembled update: the restriction is
first-order and holds per batch, it says nothing about the scores a given parameter vector
produces, and it has no effect at inference whatsoever.

The opposition should not be pressed further than it goes, and one qualification is
due on HCC's side. Because the projection sits inside the computational graph, its
Jacobian filters the backward pass as well: by \cref{prop:hcc_backward}, the
level-$(l+1)$ signal reaching the raw logits is mean-centred within each sibling
group. HCC therefore does act in gradient space, but only as a consequence of what
it does in output space, and only on the level heads' own logits rather than on the
shared parameters. It is an output-space method with a gradient-space side effect,
not a second gradient-space method; the distinction matters because the side effect
is not something the design chose, and it cannot be varied independently of the
constraint.

Whether the two can be combined is a question about attribution rather than about mechanics, and
it is settled in \S\ref{ssec:design_axis}. The full comparison of the three mechanisms, including
the LH-projection of \S\ref{sec:methodology_lhproj}, is given once in
Table~\ref{tab:framework_mechanisms}.


\section{LH-projection}
\label{sec:methodology_lhproj}

The third mechanism considered in this thesis originates in the lexicographic
hybrid deep neural network of \cite{fiaschi2024informed}, and is referred to here
as the \emph{LH-projection}. Where HCC constrains the logits emitted by the heads
(\S\ref{sec:methodology_hcc}) and lexicographic mode constrains the update
direction assembled from three separate backward passes
(\S\ref{sec:methodology_lex}), the LH-projection operates at the branch points of
a shared trunk: at the point where each level head reads the shared
representation, the component lying in the subspace already claimed by the
higher-priority heads is removed. As \S\ref{ssec:lhproj_forward} establishes, the
removal is realised in the backward pass alone, so what the mechanism filters is
the learning signal and not the representation itself.

This section derives the operator and its properties
(\S\ref{ssec:lhproj_construction}--\S\ref{ssec:lhproj_forward}), states the
instantiation used here on top of the Hier-COS taxonomy-frame architecture,
including the structural conditions that instantiation requires
(\S\ref{ssec:lhproj_hiercos}), and delimits the guarantee
(\S\ref{ssec:lhproj_limits}).

\subsection{The branch projection}
\label{ssec:lhproj_construction}

Consider a classifier with a shared trunk producing a pre-activation
$k\in\mathbb{R}^{D}$ and a representation $z=\rho(k)\in\mathbb{R}^{D}$, followed by
$L$ linear level heads with weight matrices $W_l\in\mathbb{R}^{C_l\times D}$. In
the unconstrained model every head reads the same $z$, so the level objectives
compete for the same directions of $\mathbb{R}^{D}$ with no ordering between them.

The LH-projection introduces an ordering by treating as \emph{reserved} the
directions of $\mathbb{R}^{D}$ along which the higher-priority heads' scores
change. For level $l>1$ and sample $b$ of the batch, define the stacked matrix
\[
A^{(l)}[b]
:=
\begin{bmatrix} W_1\\ \vdots\\ W_{l-1}\end{bmatrix}
\operatorname{diag}\!\big(\rho'(k_b)\big)
\in\mathbb{R}^{R_l\times D},
\qquad
R_l:=\sum_{j<l}C_j ,
\]
and let
\[
c^{(l)}
:=
\big(A^{(l)}[b]\big)^{\!\top}
\Big(A^{(l)}[b]\big(A^{(l)}[b]\big)^{\!\top}+\varepsilon I\Big)^{-1}
A^{(l)}[b]\, z
\;=:\;
P^{(l)}_\varepsilon[b]\; z .
\]
The representation nominally offered to head $l$ is $z-c^{(l)}$: the higher levels
keep the subspace they read, and level $l$ is confined to what remains. The
qualification matters, and \S\ref{ssec:lhproj_forward} discharges it: the
implementation restores the removed component in the forward pass, so the removal
takes effect on the gradient only. Both the head weights $W_j$ and the activation
derivative $\rho'(k_b)$ are taken \emph{detached} from the computational graph, so
that the reserved subspace is treated as given rather than as something the lower
level may reshape --- the parameter-space counterpart of the stop-gradient of
\S\ref{ssec:hcc_extensions}. The regularisation $\varepsilon I$ plays the same role
as there: it makes the linear system solvable when the stacked rows are close to
linearly dependent, at the price of an exactness that
\cref{prop:lhproj_operator} quantifies. The system is solved once per sample, at a
cost of a batch of $R_l\times R_l$ solves per level.

\paragraph{Why the activation derivative belongs in $A^{(l)}$.} The factor
$\operatorname{diag}(\rho'(k_b))$ is not a refinement of the operator but a
condition for it to remove the right subspace, and this is worth establishing
before anything is built on top of it. Head $j<l$ emits $W_j\rho(k)$, so the map
from the branch point to the stacked higher-level scores is
$F(k):=W_{<l}\,\rho(k)$ with Jacobian
\[
\frac{\partial F}{\partial k}\Big|_{k_b}
= W_{<l}\operatorname{diag}\!\big(\rho'(k_b)\big)
= A^{(l)}[b] .
\]
The directions in which a move at the branch point changes what the higher levels
predict are therefore the rows of $A^{(l)}[b]$, and they coincide with the rows of
$W_{<l}$ only when $\rho$ is the identity. The following proposition makes the
consequence precise.

\begin{proposition}
\label{prop:lhproj_rho}
Write $D_b:=\operatorname{diag}(\rho'(k_b))$ and let $\tilde g$ be the level-$l$
gradient with respect to $z$ after filtering, so that the level-$l$ signal at the
branch point is $\nabla_k\ell_l=D_b\,\tilde g$. Treating $k$ as the free variable,
the first-order change of the higher-level scores induced by a step along
$-\nabla_k\ell_l$ is proportional to $-W_{<l}D_b^{2}\,\tilde g$. Suppose
$D_b^{2}=D_b$, as holds for $\rho=\mathrm{ReLU}$, where $\rho'$ is a $0/1$ mask.
Then this change vanishes for every $\tilde g$ orthogonal to
$\operatorname{row}\big(W_{<l}D_b\big)$, but not, in general, for $\tilde g$
orthogonal to $\operatorname{row}(W_{<l})$.
\end{proposition}
\begin{proof}
Expanding $F$ about $k_b$ gives
$F(k_b+\delta k)-F(k_b)=W_{<l}D_b\,\delta k+O(\lVert\delta k\rVert^2)$. With $k$ as
the free variable the step is $\delta k=-\eta\,\nabla_k\ell_l=-\eta D_b\tilde g$,
which yields the stated proportionality. If $D_b^{2}=D_b$ then
$W_{<l}D_b^{2}\tilde g=W_{<l}D_b\tilde g=A^{(l)}[b]\,\tilde g$, which is zero
exactly when $\tilde g\perp\operatorname{row}(A^{(l)}[b])$. If instead $\tilde g$
is only known to satisfy $W_{<l}\tilde g=0$, no conclusion follows about
$W_{<l}D_b\tilde g$, since $D_b\tilde g$ need not lie in $\ker W_{<l}$.
\end{proof}

The batch-shared operator obtained by omitting $\rho'$ --- that is, by reserving
$\operatorname{row}(W_{<l})$ --- therefore fails to deliver the first-order
guarantee even in the ReLU case for which the guarantee was originally stated. It
removes a subspace that is not the one the higher heads are sensitive to. That form
is consequently not used anywhere in this thesis, and $A^{(l)}[b]$ above is the
operator throughout. The price is that the reserved subspace becomes
input-dependent and the linear system must be solved per sample rather than once
per batch.

In LH-DNN as originally published, $\rho$ is a ReLU and
$\rho'(k)=\mathbb{1}[k>0]$; in the Hier-COS instantiation of
\S\ref{ssec:lhproj_hiercos} a shared PReLU is used and
$\rho'(k)=\mathbb{1}[k>0]+a\,\mathbb{1}[k\le 0]$ with $a$ the learned negative
slope, taken detached. Since $D_b^2=D_b$ requires each diagonal entry to be $0$ or
$1$, the identity in \cref{prop:lhproj_rho} is exact for ReLU and only approximate
for PReLU, by an amount governed by the distance of $a$ from $\{0,1\}$; this is
recorded in \S\ref{ssec:lhproj_limits}. A second, more benign consequence of the
factorisation is that $D_b$ is invertible whenever $a\neq 0$, so
$\operatorname{rank}A^{(l)}[b]=\operatorname{rank}W_{<l}$ and the rank condition of
\cref{prop:lhproj_rank} is unaffected by the activation. Under a ReLU it would not
be: there $D_b$ is singular, and the condition would have to hold among the units
active for the given sample.

\begin{proposition}
\label{prop:lhproj_operator}
Let $A\in\mathbb{R}^{R\times D}$ have singular values $\sigma_1,\dots,\sigma_r>0$
with right singular vectors $v_1,\dots,v_r$, and let
$P_\varepsilon:=A^{\top}(AA^{\top}+\varepsilon I)^{-1}A$ with $\varepsilon>0$.
Then $P_\varepsilon$ is symmetric positive semi-definite with
\[
P_\varepsilon=\sum_{i=1}^{r}\frac{\sigma_i^2}{\sigma_i^2+\varepsilon}\,v_iv_i^{\top},
\]
so that its eigenvalues are $\sigma_i^2/(\sigma_i^2+\varepsilon)\in(0,1)$ on
$\operatorname{row}(A)$ and $0$ on $\operatorname{row}(A)^{\perp}$. Consequently
$P_\varepsilon\to P_0$, the orthogonal projector onto $\operatorname{row}(A)$, as
$\varepsilon\to 0^{+}$, and for $\varepsilon>0$ the residual retained by
$I-P_\varepsilon$ along $v_i$ is $\varepsilon/(\sigma_i^2+\varepsilon)$.
\end{proposition}
\begin{proof}
Write the thin SVD $A=U\Sigma V^{\top}$ with $U^{\top}U=V^{\top}V=I_r$ and
$\Sigma=\operatorname{diag}(\sigma_1,\dots,\sigma_r)$. Then
$AA^{\top}=U\Sigma^{2}U^{\top}$ and
\[
(AA^{\top}+\varepsilon I)^{-1}
=U(\Sigma^{2}+\varepsilon I)^{-1}U^{\top}+\tfrac{1}{\varepsilon}\big(I-UU^{\top}\big).
\]
Since $(I-UU^{\top})U=0$, the second term is annihilated on both sides by
$A^{\top}=V\Sigma U^{\top}$ and $A=U\Sigma V^{\top}$, leaving
$P_\varepsilon=V\Sigma(\Sigma^{2}+\varepsilon I)^{-1}\Sigma V^{\top}
=V\operatorname{diag}\big(\sigma_i^2/(\sigma_i^2+\varepsilon)\big)V^{\top}$,
which is the stated spectral decomposition. The remaining claims follow by
inspection of the eigenvalues.
\end{proof}

Two consequences are used later. First, the removal is exact only in the limit
$\varepsilon\to0$: for $\varepsilon>0$ a fraction
$\varepsilon/(\sigma_i^2+\varepsilon)$ of the reserved component survives, and the
fraction is largest along the \emph{weakest} singular directions of the stacked
head matrix --- that is, precisely where the higher-priority heads are least
confident of their geometry. Second, $I-P_\varepsilon$ is never rank-deficient for
$\varepsilon>0$, so the operator is always well defined; whether it is
\emph{informative} is a separate question, settled by the rank condition of
\S\ref{ssec:lhproj_hiercos}.

\subsection{What the projection acts on}
\label{ssec:lhproj_forward}

What head $l$ receives is not $z-c^{(l)}$ but
\[
z^{(l)} \;:=\; z-c^{(l)}+\operatorname{sg}\big(c^{(l)}\big),
\]
and this changes the character of the mechanism in a way that must be made
explicit, because it determines whether the LH-projection alters the function
computed by a given parameter vector or only the parameters that training
reaches.

\begin{proposition}
\label{prop:lhproj_forward}
Let $c^{(l)}=P^{(l)}_\varepsilon[b]\,z$ with $P^{(l)}_\varepsilon[b]$ constant with
respect to $z$, and let
$z^{(l)}=z-c^{(l)}+\operatorname{sg}(c^{(l)})$. Then
\[
\text{(i)}\quad z^{(l)}=z
\qquad\text{(forward value)},
\qquad\qquad
\text{(ii)}\quad
\frac{\partial z^{(l)}}{\partial z}=I-P^{(l)}_\varepsilon[b]
\qquad\text{(Jacobian)} .
\]
\end{proposition}
\begin{proof}
(i) The stop-gradient operator returns its argument unchanged in the forward pass,
so $z^{(l)}=z-c^{(l)}+c^{(l)}=z$. (ii) In the backward pass
$\operatorname{sg}(c^{(l)})$ contributes no derivative, so
$\partial z^{(l)}/\partial z=I-\partial c^{(l)}/\partial z$. The matrix
$P^{(l)}_\varepsilon[b]$ is built from detached head weights and a detached
activation derivative; it is therefore constant with respect to $z$, and
$c^{(l)}=P^{(l)}_\varepsilon[b]\,z$ is linear in $z$ with
$\partial c^{(l)}/\partial z=P^{(l)}_\varepsilon[b]$.
\end{proof}

The mechanism is therefore \emph{not} a modification of the forward map. Head $l$
receives the unprojected representation $z$, and the function computed by a given
parameter vector is exactly that of the unconstrained multi-head classifier. What
the projection modifies is the gradient: by (ii), the contribution of $\ell_l$ to
the gradient reaching the shared trunk is
\[
\big(I-P^{(l)}_\varepsilon[b]\big)\,\nabla_{z}\,\ell_l ,
\]
that is, the component of the level-$l$ learning signal lying along the directions
that move the higher levels' scores is removed before it reaches the shared
parameters. The LH-projection is therefore an intervention on the update direction
and not on the forward map, and it has no effect at inference.

This places it beside lexicographic mode rather than beside HCC, and the two admit
a direct comparison, which is one reason both are studied here:

\begin{itemize}
    \item \textbf{Reference directions.} Lexicographic mode orthogonalises against
    the \emph{gradients} of the higher-priority objectives, which change at every
    step and depend on the current batch. The LH-projection orthogonalises against
    the rows of the higher heads' \emph{Jacobian at the branch point}, a geometric
    object that varies slowly through the weights and per sample through the
    activation mask. The former asks ``in which direction does the coarse loss want
    to move?''; the latter asks ``which part of the representation does the coarse
    level currently read?''.
    \item \textbf{Granularity.} The LH-projection is applied per sample, inside the
    forward graph, before any gradient is accumulated; lexicographic mode is
    applied per batch, to the accumulated gradient, on parameter blocks determined
    by gradient support (\S\ref{ssec:lex_blocks}).
    \item \textbf{Locus.} The LH-projection filters the signal at a single point,
    the branch point of the shared trunk, and therefore cannot control how the
    level objectives interact in parameters upstream of it beyond what that filter
    removes. Lexicographic mode acts directly on the parameters of every shared
    block.
    \item \textbf{Cost.} The LH-projection requires one backward pass and $L-1$
    linear solves of size $R_l$ per sample; lexicographic mode requires $L$
    backward passes and no solves.
\end{itemize}

\subsection{Instantiation on Hier-COS}
\label{ssec:lhproj_hiercos}

The LH-projection is defined for a shared trunk read by independent linear level
heads. Hier-COS as published does not have that structure: it emits a single node
vector, maps it through a \emph{fixed} frame shared by the whole taxonomy, and
scores levels by subspace norms (\S\ref{ssec:problem_scores}). Transplanting the
projection onto it therefore requires three structural changes, which are part of
the construction rather than free hyperparameters, and which are stated here for
that reason.

\paragraph{(i) Independent level heads, hence an identity frame.} The projection
is defined in terms of the weight matrices $W_1,\dots,W_{l-1}$ of the heads above
level $l$. Under the dense random orthonormal frame of the published model the
scores of every level are linear mixtures of a single shared node space, and no
per-level matrix exists whose rows could be stacked into $A^{(l)}$: the object the
construction needs is simply not present. The projected variant therefore replaces
the frame by the identity --- equivalently, by a per-level block-diagonal frame,
which leaves the level blocks independent --- and attaches one learnable linear
head per level to the transform output. The combination of the projection with a
dense frame is therefore not admitted at all. The frame is thus a
\emph{precondition} of the method, and the identity-versus-random comparison
measures the cost of the precondition, not a tuning choice; that cost is reported
in the experimental chapter.

\paragraph{(ii) Level-scoped normalisation, not merely a decomposed objective.}
The objective must expose one differentiable term per level, and it must do so in
a form that each level's branch can receive on its own. Two steps separate the
published objective from that. The published objective is a KL divergence between
one softmax taken over all $N$ taxonomy nodes and a target node distribution: it
yields a single scalar, and no per-level signal can be extracted from it at all.
An intermediate objective, retaining the global softmax but reading off the target
node of each level separately, does yield three differentiable terms and is
sufficient for the gradient-space ordering of \S\ref{sec:methodology_lex}. It is
\emph{not} sufficient here: the three terms share one normaliser over all nodes, so
the gradient of the level-$l$ term flows back through the branches of every other
level, and the filtered representation that branch $l$ reads would no longer
determine that term's learning signal. Only the level-scoped objective, in which
the softmax is taken within each level's node block, makes each level loss a
function of its own branch alone, and it is the one the projected variant uses,
together with the same node-norm regulariser. As with the frame, the comparison
against the published objective quantifies the price of a structural requirement.

\paragraph{(iii) A non-trivial complement.} The projection is informative only if
the reserved subspace does not exhaust the representation.

\begin{proposition}
\label{prop:lhproj_rank}
Let $D$ be the width of the shared representation. If $D\le R_L=\sum_{l<L}C_l$
then, for $A^{(L)}$ of full row rank, $\operatorname{row}(A^{(L)})=\mathbb{R}^{D}$
and $I-P^{(L)}_0=0$: the projection annihilates the entire learning signal of the
finest level. A non-trivial complement at every level requires
$D>\sum_{l<L}C_l$.
\end{proposition}
\begin{proof}
$A^{(L)}$ has $R_L$ rows in $\mathbb{R}^{D}$; if it has full row rank and
$R_L\ge D$ its row space is all of $\mathbb{R}^{D}$, so the orthogonal projector
onto that row space is the identity and its complement is zero. Since $R_l$ is
increasing in $l$, the binding constraint is $l=L$.
\end{proof}

The condition $D>\sum_{l<L}C_l$ is verified before training rather than left to be
discovered from the results, and it requires the width of the representation to be
decoupled from the taxonomy size $N$ that the unprojected model uses. The reason
for checking it in advance is that for $\varepsilon>0$ the failure predicted by
\cref{prop:lhproj_rank} is not a hard zero but a uniform shrinkage of
the fine-level gradient by $\varepsilon/(\sigma_i^2+\varepsilon)$
(\cref{prop:lhproj_operator}), which would be hard to distinguish from
a badly scaled learning rate.

\paragraph{Remaining degrees of freedom.} Given the three requirements above, the
following are genuine ablation axes, varied in the experimental chapter and not
constitutive of the method: the placement of the projection relative to the
transformation layer, that is, whether the projected representation is the full
residual transform, a batch-norm-plus-linear reduction, or the backbone output
itself; the representation width, subject to \cref{prop:lhproj_rank}; and the
\emph{advantage} parameterisation, in which the logits of level $l$ are expressed
as a residual on top of the detached logit of the predicted parent,
$z_l \leftarrow z_l + z_{l-1}[\pi_{l-1}(\cdot)]$, so that the head predicts a
correction to its parent rather than an absolute score. The composition of the
reserved subspace is \emph{not} on this list: by \cref{prop:lhproj_rho} the
activation Jacobian is part of the operator, not a variant of it.

\paragraph{Exclusivity with HCC.} The LH-projection and HCC are mutually exclusive
on Hier-COS: the configuration validator rejects the pair and no run combines
them. The reason is that they would act on the same quantity in incompatible ways.
The LH-projection determines the per-level logits by filtering what each head
reads, and its reserved subspace is, by \cref{prop:lhproj_rho}, the row space of
the Jacobian of the higher heads' outputs with respect to the branch point. HCC
overwrites those same logits with the solution of an affine problem anchored on the
level above. Composing them would leave the reserved subspace describing the
Jacobian of a map that no longer produces the logits the objective evaluates, so
the operator would remove a subspace with no relation to what it is meant to
protect.

\subsection{Scope of the guarantee}
\label{ssec:lhproj_limits}

The guarantee attaching to the LH-projection in \cite{fiaschi2024informed} is
narrower than the construction may suggest, and four qualifications should be
carried into the interpretation of the results. They parallel those of
\S\ref{ssec:lex_guarantee}.

\begin{enumerate}
    \item \textbf{First-order.} The statement concerns the gradient at the current
    parameters. It guarantees that the level-$l$ signal reaching the trunk carries
    no component along the reserved subspace at that point; it does not guarantee
    that the higher-priority objectives do not increase after the step, and says
    nothing about the limit of the trajectory. As with HCC and with lexicographic
    mode, this is lexicographic \emph{structure} without lexicographic
    \emph{sequential optimality}.
    \item \textbf{Per-sample and branch-point-local.} The projection is applied at
    the branch point, on the representation each head reads. Interactions between
    level objectives that are mediated by parameters upstream of the branch point,
    or by batch statistics such as those of the normalisation layers, are outside
    its scope.
    \item \textbf{Conditional on the activation being idempotent in its
    derivative.} By \cref{prop:lhproj_rho} the guarantee closes only when
    $\operatorname{diag}(\rho'(k_b))$ is idempotent, that is, when every entry of
    $\rho'$ is $0$ or $1$. This is exact for the ReLU of the published
    formulation. The Hier-COS instantiation uses a shared PReLU with a learned
    negative slope $a$, taken detached, for which the condition fails unless $a$
    happens to reach $0$ or $1$; the guarantee is then approximate, with an error
    governed by $a$ and by the fraction of coordinates in the negative branch.
    The slope is a learned parameter, so this is a property to be checked on the
    trained model rather than assumed, and it is one of the reasons a projected
    Hier-COS run is not simply LH-DNN with a different backbone.
    \item \textbf{Approximate, and least accurate where it matters most.} By
    \cref{prop:lhproj_operator} the removal leaves a residual
    $\varepsilon/(\sigma_i^2+\varepsilon)$ along the $i$-th singular direction of
    $A^{(l)}[b]$. The default $\varepsilon=10^{-6}$ makes this negligible for
    well-conditioned heads, but the residual grows as $\sigma_i\to 0$, that is,
    early in training and along class directions the higher levels have not yet
    separated.
\end{enumerate}


\section{Hierarchical inference strategies}
\label{sec:methodology_inference}

The preceding sections described how a network is trained. Turning the trained
network into a predicted path is a separate decision, and one that materially
changes every reported number: the same checkpoint read out and decoded in two
different ways yields different accuracies, a different full-path accuracy and,
in particular, a different measured inconsistency.

The decision factorises into two independent choices, and the section is
organised accordingly:

\begin{enumerate}
    \item a \textbf{readout} $u\mapsto(z_1,\dots,z_L)$, which reduces the
    network's internal taxonomy-node coordinates to one score vector per level
    (\S\ref{ssec:inference_readout});
    \item a \textbf{decoder} $(z_1,\dots,z_L)\mapsto\hat y$, which turns level
    scores into a path (\S\ref{ssec:inference_decoding}).
\end{enumerate}

Both are properties of the evaluation, not of the trained model. Neither touches
the parameters, so both can be replaced after training on a frozen checkpoint ---
which is what makes the grid of \S\ref{ssec:inference_grid} an instrument rather
than a convention. The section then states the checkpoint-selection rule that
keeps decoder comparisons honest (\S\ref{ssec:inference_selection}), and closes by
recording exactly which of the training-time mechanisms can and cannot move which
reported quantity (\S\ref{ssec:inference_interaction}). The metrics computed on
top of a chosen cell of the grid are defined with the experimental setup, in
\cref{ssec:inference_metrics}, and are referred to here by name.

\subsection{Readout rules}
\label{ssec:inference_readout}

Every family considered here can be viewed as emitting a vector of
\emph{node coordinates} $u\in\mathbb{R}^{N}$ over the taxonomy, partitioned into
per-level blocks. For the classifier-head families (H-CAST, HRN, HT-CapsNet,
LH-DNN) the blocks are simply the per-level logits of
\S\ref{ssec:problem_scores}; for Hier-COS they are the coordinates of the frame
output. Two readouts are defined on this common representation.

\paragraph{Node score.} Each node is ranked by its own coordinate,
\[
z_l[c] \;=\; \phi\big(u[c]\big),
\]
with $\phi$ the identity for signed classifier logits and $\phi=\lvert\cdot\rvert$
for Hier-COS frame coordinates. The distinction is not cosmetic. A classifier
logit is signed and a large negative value carries the meaning \emph{not this
class}; Hier-COS, by contrast, trains its node distribution as
$\operatorname{log\,softmax}(\lvert u\rvert)$, so its objective never constrains
the sign of a coordinate and only the magnitude carries class evidence. In both
cases the rule ranks nodes by the quantity that the model's own objective drives
up for the correct node.

\paragraph{Subspace norm.} Each node is ranked by the $L_2$ norm of the projection
of $u$ onto the taxonomy subspace $S_c$ assigned to that node, comprising its
ancestors, itself and its descendants:
\[
z_l[c] \;=\; \big\lVert \Pi_{S_c} u \big\rVert_2 .
\]
This is the native rule of Hier-COS (\S\ref{ssec:problem_scores}). Applied to a
classifier-head model it constitutes an \emph{inference-only assumption}: the
per-level logits are treated as coordinates in an identity taxonomy frame. The
assumption is worth stating explicitly because it is not innocuous --- a norm
squares its inputs, so the subspace-norm readout discards the sign of a signed
classifier logit unconditionally. That loss of the ``not this class'' evidence is
the substantive content of the identity-frame assumption, and it is the reason
this readout is expected to behave differently on the two kinds of family.

\paragraph{Each model's native rule is one readout, not a separate case.} The
readouts above are defined for every family, and each family's own inference
procedure is recovered as one of them: \texttt{node\_score} for the
classifier-head families, \texttt{subspace\_norm} for native Hier-COS.
Reproduction is exact for H-CAST and Hier-COS, which recompute the same quantity,
and metric-exact for HRN and the remaining classifier heads, where the readout
reads raw logits rather than the sigmoid or softmax the model applies: those maps
are monotone within a level, and every metric of \S\ref{ssec:inference_metrics} is
computed from an argmax, so no reported quantity changes. Treating the native rule
as a grid cell rather than a privileged baseline is what allows a model to be
scored under another family's readout without special-casing.

\subsection{Decoding rules}
\label{ssec:inference_decoding}

\paragraph{Independent decoding.} Each level is decoded on its own,
\[
\hat y_l=\arg\max_{c}\; z_l[c],\qquad l=1,\dots,L .
\]
No structural information is used. The resulting path need not be consistent: the
predicted fine class may belong to a different parent than the predicted coarse
class. Independent decoding therefore measures what the model has learned at each
level in isolation, and its inconsistency rate is a genuine diagnostic of whether
the training-time mechanisms produced hierarchically coherent scores.

\paragraph{Top-down decoding.} The coarsest level is decoded first and each
subsequent level is restricted to the children of the level above:
\[
\hat y_1=\arg\max_{c}\;z_1[c],
\qquad
\hat y_{l}=\arg\max_{c\,:\,\pi_{l-1}(c)=\hat y_{l-1}}\; z_{l}[c],
\quad l=2,\dots,L .
\]
Should a predicted parent have no children --- which the ``no childless parent''
assumption of \S\ref{sec:methodology_problem} excludes, but which is guarded
against nonetheless --- the level falls back to an unrestricted argmax for the
affected samples.

\begin{proposition}
\label{prop:topdown_consistency}
If the taxonomy satisfies the ``no childless parent'' assumption, the path
produced by top-down decoding is consistent for every input.
\end{proposition}
\begin{proof}
By induction. $\hat y_1$ is unconstrained. For $l>1$, the feasible set
$\{c:\pi_{l-1}(c)=\hat y_{l-1}\}$ is non-empty by assumption, so the masked argmax
is attained on it and $\pi_{l-1}(\hat y_l)=\hat y_{l-1}$ holds by construction; the
fallback branch is therefore never taken.
\end{proof}

The two decoders are not ordered by quality, and the trade-off between them is one
of the objects of study of this thesis. Top-down decoding purchases consistency at
the price of \emph{error propagation}: a mistake at a coarse level removes the
correct fine class from the feasible set, so that a sample which independent
decoding would have classified correctly at the fine level is necessarily
misclassified. Independent decoding avoids this coupling but offers no structural
guarantee. Consequently top-down decoding cannot change the level-$1$ prediction,
can only alter levels $l\ge2$, and may lower fine-level accuracy relative to
independent decoding while lowering the inconsistency rate to zero.

\subsection{The inference grid and its post-hoc use}
\label{ssec:inference_grid}

The readout of \S\ref{ssec:inference_readout} may be preceded by a
\emph{transform} applied to the node coordinates. One transform is considered: the
HCC affine hierarchy projection of \S\ref{sec:methodology_hcc}, applied to $u$
before the readout consumes it. Crossing the two readouts with the presence or
absence of that transform gives four inference rules,
\[
\big\{\texttt{node\_score},\ \texttt{subspace\_norm}\big\}
\times
\big\{\text{none},\ \text{HCC}\big\},
\]
and crossing those with the two decoders of \S\ref{ssec:inference_decoding} gives
the eight cells under which a checkpoint can be scored. Every cell is defined for
every family --- no combination is rejected --- and all of them are computed from
a single shared forward pass, since they differ only in what is done to $u$
afterwards.

\paragraph{Why this is an instrument and not a convention.} Nothing in the grid
trains or modifies a parameter: it is applied to an already-selected checkpoint,
with no loss, optimiser or training loop constructed. That is precisely what makes
it useful. A training-time mechanism that improves the reported metrics admits two
readings --- it changed the representation, or it supplied a better rule for
reading an unchanged representation --- and the grid separates them by
construction. Applying the HCC transform post-hoc to a checkpoint trained
\emph{without} HCC isolates what the constraint is worth as a pure output
correction; comparing that against a checkpoint trained \emph{with} HCC isolates
what the constraint is worth as a training signal. The difference between the two
is the part of HCC's effect that is genuinely representational, and it cannot be
obtained from either run alone. The same logic applies to the frame question of
\S\ref{ssec:design_substrate}: whether Hier-COS's advantage lives in its
subspace-norm readout can be tested by giving that readout to a model that was
never trained with one.

Two properties of the grid follow from the rules themselves, and should not be
mistaken for empirical findings when the results are read.

\begin{proposition}
\label{prop:posthoc_hcc_topdown}
Let the readout be \texttt{node\_score} with $\phi$ the identity, and let the
transform be the HCC projection. Then the transform cannot change any metric
computed under top-down decoding.
\end{proposition}
\begin{proof}
By \S\ref{ssec:hardnet_to_hcc} the correction is
$\hat z_{l+1} = z_{l+1} - M^{\top}(MM^{\top}+\varepsilon I)^{-1}(Mz_{l+1}-\hat z_l)$.
By \cref{lemma:diag} the inverted matrix is diagonal, so the vector
subtracted from $z_{l+1}$ assigns to every child the same value, namely the
coefficient of its parent; within the children of a fixed parent the correction is
therefore a common additive shift, which preserves their ordering. Top-down
decoding evaluates an argmax restricted to exactly such a sibling set
(\S\ref{ssec:inference_decoding}). The coarsest level is anchored, $\hat z_1=z_1$,
so $\hat y_1$ is unchanged; inductively, each subsequent masked argmax is taken
over a set whose internal ordering the correction preserves, so $\hat y_l$ is
unchanged for every $l$.
\end{proof}

The argument establishes slightly more than the statement, and the surplus is
worth extracting, because it identifies the only route by which HCC can change an
independently decoded prediction at all.

\begin{corollary}
\label{cor:hcc_sibling_shift}
Under the hypotheses of \cref{prop:posthoc_hcc_topdown}, the HCC
correction adds to every child of a fixed parent the same scalar. Consequently it
never reorders two classes sharing a parent, and any change it produces in an
independently decoded prediction, and hence in
$\mathrm{TICE}_{\text{independent}}$, arises exclusively from the reordering of
classes belonging to \emph{different} parents.
\end{corollary}
\begin{proof}
By \cref{lemma:diag} the matrix inverted in the correction is diagonal, so
the vector $M^{\top}(MM^{\top}+\varepsilon I)^{-1}(Mz_{l+1}-\hat z_l)$ has, in the
coordinate of child $j$, the value of the coefficient of its parent $\pi_l(j)$
alone. Two children of the same parent therefore receive an identical additive
term, which cancels in their difference and preserves their relative order.
Independent decoding is an unrestricted argmax over the whole level, so the only
comparisons whose outcome the correction can alter are those between classes with
distinct parents, whose additive terms differ.
\end{proof}

This is the mechanistic content of the claim, made in
\S\ref{ssec:inference_interaction}, that HCC should be expected to reduce rather
than to eliminate the independent-decoding inconsistency: the correction can move
probability mass between sibling groups, which is what a consistency gain
requires, but it is structurally incapable of resolving a competition inside one.

The proposition fails for the magnitude convention $\phi=\lvert\cdot\rvert$ used
by Hier-COS, where a common additive shift is not order-preserving, and it fails
for the subspace-norm readout, whose subspaces span several levels: there the
transform moves both decoders for every family. The practical consequence is a
check rather than a result --- a post-hoc HCC row differing from its
untransformed counterpart under top-down decoding on a classifier-head model
indicates a defect in the evaluation, not an effect.

\paragraph{Three qualifications on the post-hoc HCC cells.} First, the transform is
a faithful reproduction of the training-time projection only in that it applies the
same operator to the same coordinates; it reuses the taxonomy of the run under
evaluation, and any mismatch between that taxonomy and the one the checkpoint was
trained on would surface here rather than being caught elsewhere.
Second, the transform acts on whatever
coordinates the readout consumes, which for HRN is its auxiliary leaf
cross-entropy head rather than the sigmoid tree branch that HCC corrects during
training; the HRN cells are therefore not a reconstruction of that family's
training-time HCC path and are not reported as one. Third, the grid inherits the
hazards of any re-evaluation of an already completed run: if the test split a run
was scored on cannot be reconstructed exactly and a substitute is used, the
rebuilt label space may differ from the one the checkpoint trained on, which
surfaces as plausible but meaningless parent-level metrics while leaf accuracy
still matches. Every reported cell is therefore checked against the run's own
test metrics on its native cell before any other row from that run is used.

\subsection{Checkpoint selection}
\label{ssec:inference_selection}

Because the two decoders induce different orderings of the epochs, a single
``best'' checkpoint would silently favour whichever decoder the selection metric
happened to reward. Every run therefore maintains \emph{two} checkpoints, one per
decoding mode, each selected on the validation split by the mode's own metrics.

The selection criterion is itself lexicographic, which is convenient given the
subject of this thesis. Epochs are ranked by the tuple
\[
\Big(\;\mathrm{FPA}_{\text{mode}},\;\; -\mathrm{TICE}_{\text{mode}},\;\;
\mathrm{wAP}_{\text{mode}}\;\Big)
\]
under the usual lexicographic order on tuples, so that full-path accuracy
dominates, consistency breaks ties, and weighted accuracy breaks any remaining
tie. Missing or non-finite entries rank below every finite value; when none of the
three hierarchy metrics is available the deepest-level accuracy is used as the
sole criterion.

The reporting rule that follows is applied without exception in the experimental
chapters: \textbf{a top-down result is read from the top-down-selected checkpoint
and an independent result from the independent-selected checkpoint}. The two rows
of a table may therefore come from different epochs of the same run, and mixing
them --- reporting, for instance, the independent-decoded metrics of the
top-down-selected checkpoint --- would confound the decoder comparison with a
checkpoint comparison.

\subsection{Interaction with the training-time mechanisms}
\label{ssec:inference_interaction}

The results established above, read together with the derivation of each of the
three training-time mechanisms, yield a clean division, which is worth stating
because it is easy to assume otherwise.

\begin{itemize}
    \item \textbf{Lexicographic mode and the LH-projection have no effect at
    inference.} The first modifies the update direction by construction; the
    second does so by \cref{prop:lhproj_forward}. Both change the
    parameters reached by training, and neither changes the function computed by a
    given parameter vector. Any difference they produce in the metrics above is
    entirely mediated by the learned parameters.
    \item \textbf{HCC does act at inference}, since it modifies the logits
    themselves and the same projection is applied in evaluation as in training.
    \item \textbf{HCC does not by itself guarantee consistency.} This is the point
    most easily misread. The cascade enforces $M_{l,l+1}\hat z_{l+1}=\hat z_l$, an
    equality in \emph{logit space} between the parent logit and the sum of its
    children's logits, as emphasised in \S\ref{ssec:hardnet_to_hcc}. It does not
    imply
    $\pi_l\big(\arg\max_j \hat z_{l+1}[j]\big)=\arg\max_i \hat z_l[i]$: a parent
    may hold the largest logit while the largest child logit overall belongs to a
    different parent whose children sum to less. \cref{cor:hcc_sibling_shift}
    sharpens this: since the correction shifts every sibling group by a single
    scalar, it can only reorder classes across parents, so whatever consistency it
    buys is bought there and nowhere else. HCC should therefore be expected
    to \emph{reduce} $\mathrm{TICE}_{\text{independent}}$, not to zero it, and any
    claim that it produces consistent predictions must be supported by the measured
    independent-decoding TICE rather than deduced from the constraint. Exact
    consistency at inference is delivered by top-down decoding alone.
\end{itemize}

The division has a methodological consequence that the locus questions of
\S\ref{ssec:design_locus} rely on. Because lexicographic mode and the
LH-projection cannot act at inference, any effect they produce is necessarily
representational, and no post-hoc analysis can either reproduce or explain it
away. Because HCC \emph{can} act at inference, the same is not true of it: an
improvement under HCC is a priori compatible with a purely corrective reading,
and separating the two requires the post-hoc grid of
\S\ref{ssec:inference_grid}. The asymmetry is not a nuisance but the reason the
three mechanisms cannot be compared on their headline numbers alone.




\section{The mechanism axis and its admissible configurations}
\label{sec:methodology_axis}

The chapter has so far introduced the three mechanisms one at a time, each in the
terms of its own construction, and has deferred every comparison between them.
This section collects them. It asks what the three are \emph{relative to each
other}, and which of their combinations exist at all --- questions that belong
here, with the derivations, because they are settled by the structure of the
mechanisms rather than by any measurement.

The answers turn out to be more restrictive than the three independent
configuration switches suggest, and the restriction is itself a result. The
section first classifies the mechanisms by the point of the computation at which
they act (\S\ref{ssec:design_mechanisms}); it then shows that no two of them are
used together, so that the mechanism choice is a four-way alternative rather than
a factorial, and distinguishes carefully between the two pairs that cannot be
composed and the one pair that could be but is not
(\S\ref{ssec:design_axis}); it records which families admit which mechanism, and
why the exclusions are arguments rather than gaps (\S\ref{ssec:design_support});
and it closes with the structural constraints that cut the nominal configuration
space to a tractable one (\S\ref{ssec:design_space}).

What this section does \emph{not} do is state the experiments. The research
questions, the falsification arms, the controlled variables and the reporting
protocol depend on the datasets, metrics and seed counts that
\cref{ch:experiments} introduces, and are stated there, in
\S\ref{sec:results_design}, once those are available.

\subsection{Where each mechanism acts}
\label{ssec:design_mechanisms}

\begin{table}[htbp]
\centering
\small
\begin{tabular}{@{}lllll@{}}
\toprule
\textbf{Intervention} & \textbf{Acts on} & \textbf{Granularity} & \textbf{Guarantee} & \textbf{At inference} \\
\midrule
HCC & logits & per sample & exact (affine equality) & yes \\
LH-projection & update direction & per sample & first order, at the branch point & no \\
Lexicographic mode & update direction & per batch & first order, per block & no \\
\midrule
Readout rule & level scores & per sample & --- & yes, only \\
Top-down decoding & decoded path & per sample & exact (path consistency) & yes, only \\
\bottomrule
\end{tabular}
\caption{The interventions, classified by the point of the computation at which
they act. The upper block contains the three priority mechanisms, the lower block
the two components of the inference rule. Only HCC and the inference rules change
the prediction produced by a fixed parameter vector; the two gradient-space
mechanisms act solely through the parameters that training reaches.}
\label{tab:framework_mechanisms}
\end{table}

The three priority mechanisms form the axis announced in
\S\ref{ssec:lexicographic}: priority imposed on the output contract by
construction (HCC), priority imposed on the learning signal at a structural
bottleneck (the LH-projection), and priority imposed on the assembled update in
parameter space (lexicographic mode). None of them achieves sequential
lexicographic optimality; all three realise lexicographic \emph{structure} at a
different locus. Whether the locus matters is the question the design exists to
answer.

The lower block of the table is what makes that question answerable. Because the
readout and the decoder can be varied on a frozen checkpoint
(\S\ref{ssec:inference_grid}), the contribution of a mechanism to the learned
representation can be distinguished from its contribution to the readout, which
for HCC --- the one mechanism active at inference --- is not otherwise separable.

\subsection{The mechanism axis is a single choice, not a factorial}
\label{ssec:design_axis}

It is natural to expect three independently configurable mechanisms to yield a
$2\times2\times2$ design. They do not: no two of the three are ever used together.
The three exclusions are not, however, of the same kind, and conflating them would
overstate what the design establishes. Two of the pairs cannot be composed at all;
the third could be, and is excluded by decision.

\begin{itemize}
    \item \textbf{HCC and the LH-projection} both determine the per-level logits,
    by incompatible routes: the LH-projection fixes what each head reads, while
    HCC overwrites what each head emits with the solution of an affine problem
    anchored on the level above. Under HCC the head weights whose row space the
    LH-projection reserves no longer describe the map that produces the logits the
    objective sees, so the reserved subspace would be built from a Jacobian that is
    no longer the Jacobian of anything the objective evaluates. The pair is
    rejected at construction, for the reason given in
    \S\ref{ssec:lhproj_hiercos}, and the configuration validator refuses it.
    \item \textbf{The LH-projection and lexicographic mode} are, by the analysis
    of \S\ref{ssec:lhproj_forward}, two realisations of one principle: both
    remove from a lower-priority level's learning signal the component lying in a
    subspace claimed by the higher levels, one inside the computational graph and
    one on the assembled gradient. Composing them would apply the same principle
    twice by two routes, and the resulting ordering would not be attributable to
    either. The pair is likewise rejected at construction, and the configuration
    validator refuses it.
    \item \textbf{HCC and lexicographic mode} \emph{can} be composed, and the
    reason is worth stating precisely, because the obvious argument against the
    combination does not survive inspection. One might expect HCC to destroy the
    separability of the three level objectives that \S\ref{ssec:lex_setting}
    presupposes, since each $\hat z_l$ depends on the raw logits of every level at
    or above $l$. It does not, and the stop-gradient is exactly why: each anchor is
    detached, so $\nabla_\theta\,\ell_l(\hat z_l)$ still flows through the level-$l$
    path alone, and three independently computable gradients remain available for
    the ordering to act upon. Nothing in either construction breaks, and the
    configuration validator permits the pair.

    The objection is instead one of attribution. Under HCC the coarse-to-fine
    ordering is already imposed in the forward pass and the reverse influence
    already blocked in the backward pass; adding lexicographic mode would impose
    the same ordering a second time, on gradients of objectives that HCC has made
    hierarchically dependent in their arguments. \Cref{prop:lex_firstorder} would
    continue to hold as algebra, but its reading would change: the quantity it
    protects would no longer be the level-$l$ objective but the level-$l$ objective
    \emph{as evaluated on constrained logits}, and no observed effect could be
    assigned to one locus rather than the other. The combination is therefore not
    run in this thesis, and this is a design decision rather than an impossibility.
\end{itemize}

The mechanism axis is therefore a four-way single choice,
\[
\big\{\,\text{none},\ \text{HCC},\ \text{LH-projection},\ \text{lexicographic mode}\,\big\},
\]
and every comparison among the three is a between-arms comparison against the
shared \emph{none} baseline.

Two remarks are due. First, the two genuine incompatibilities are not a limitation
of the design but one of its results, and they share a cause: both concern what
determines the per-level logits. HCC and the LH-projection both claim that role and
give incompatible answers; the LH-projection and lexicographic mode give the same
answer twice. Second, and by contrast, the HCC--lexicographic exclusion rests on
what was actually run rather than on what could be. No run in this thesis combines
them, and a study of the combination --- which the framework would support without
modification --- is left as an avenue this design deliberately does not take.

\subsection{Which families support which mechanism}
\label{ssec:design_support}

\begin{table}[htbp]
\centering
\small
\begin{tabular}{@{}lccc l@{}}
\toprule
\textbf{Family} & \textbf{HCC} & \textbf{Lex.\ mode} & \textbf{LH-proj.} & \textbf{Condition} \\
\midrule
H-CAST   & yes & yes & no  & three levels; no global KL term for lex.\ mode \\
Hier-COS & yes & yes & yes & decomposed loss for lex.; level softmax and identity frame for LH-proj. \\
HRN      & yes & yes & no  & exactly three levels; decomposed level-marginal loss \\
HT-CapsNet & yes & yes & no & three capsule margin losses \\
LH-DNN   & no  & no  & built in & no three independent level objectives \\
\bottomrule
\end{tabular}
\caption{Support matrix. ``Built in'' indicates that the mechanism is constitutive
of the family rather than optional. ``Yes'' records that a family admits a
mechanism, not that the mechanisms may be combined: by \S\ref{ssec:design_axis}
each run carries at most one.}
\label{tab:framework_support}
\end{table}

Three entries deserve comment, and the first is the most consequential for the
design.

\paragraph{LH-DNN carries no mechanism arms, and this is an argument rather than a
gap.} Every mechanism in this thesis acts on the part of the network where a
shared representation is distributed to level-specific heads --- which is
precisely the part that constitutes LH-DNN's contribution. Enabling lexicographic
mode on it would apply, by explicit gradient surgery, the same principle its
branch-point projection already applies inside the graph; imposing HCC on it would
overwrite the structure it exists to demonstrate. Any arm applied to LH-DNN
therefore either duplicates what it already does or replaces it, and a result from
such an arm could not be attributed to LH-DNN. It is retained as a baseline and as
the origin of the LH-projection, and it is the right control precisely because it
admits no arms: it holds the position ``hierarchy handled natively by the
architecture'' fixed while the other four families vary.

\paragraph{Hier-COS is the only family supporting all three mechanisms}, which is
why it carries the structural ablations of \S\ref{ssec:lhproj_hiercos} and serves
as the common substrate on which the loci can be compared at matched backbone,
frame, loss and weights. The head-to-head between the LH-projection and
lexicographic mode --- the two realisations of one principle --- is therefore
available \emph{within} a single family, holding everything but the mechanism
fixed, which is a cleaner contrast than a cross-family comparison could have been.

\paragraph{HRN's three-level restriction} is a property of the method as published,
not of the hierarchy, and coincides here with the three-level restriction that HCC and
lexicographic mode impose for their own reasons.

\subsection{The design space and what prunes it}
\label{ssec:design_space}

Enumerated as a factorial, the space is intractable and uninformative. On Hier-COS
alone the frame ($3$), the loss scope ($3$), the level weighting ($5$), the
placement of the learnable transform ($3$), the mechanism choice ($8$, counting
the two composition modes crossed with the two removal rules of
\S\ref{ssec:lex_projection} as distinct arms) and the dataset ($3$) already give
$3\times3\times5\times3\times8\times3 = 3240$ configurations per seed, before any
backbone or head variant, and at three seeds each. The factorial is not the plan.

What cuts it to the order of a few dozen configurations that carry information is
a set of structural constraints, each of which rules a region of the space out
rather than merely making it unattractive. They are worth stating because several
of them are themselves substantive:

\begin{enumerate}
    \item \textbf{Gradient-space ordering requires separable level objectives, and
    the two mechanisms demand separability at two different strengths.}
    Lexicographic mode requires three differentiable level losses, which rules out
    any objective returning a single coupled scalar over all taxonomy nodes ---
    Hier-COS's published KL objective, and H-CAST's global KL term. It does not
    rule out a global softmax whose per-level target terms are read off separately,
    since those are three differentiable scalars. The LH-projection asks for more:
    each level loss must be a function of its own branch alone, which a shared
    normaliser destroys, so it additionally requires the softmax to be taken within
    each level's node block (\S\ref{ssec:lhproj_hiercos}). Separability is
    therefore not one condition but a graded one, and where a family's objective
    sits on that scale determines which mechanisms it admits.
    \item \textbf{The LH-projection requires independent level heads}, hence an
    identity or per-level block-diagonal frame, and a representation wider than
    the total number of non-leaf classes (\cref{prop:lhproj_rank}).
    \item \textbf{HCC requires exactly three levels}, as do the three-objective
    constructions of lexicographic mode.
    \item \textbf{Family-specific decompositions}, recorded in
    Table~\ref{tab:framework_support}: H-CAST without its global KL term, the
    level-marginal objective on HRN, the three capsule margin losses on
    HT-CapsNet.
\end{enumerate}

The pruning is therefore not a convenience. The constraints trace the boundary of
where each mechanism is definable at all, and that boundary is treated in
\cref{ch:experiments} as an object of study rather than as a precondition to be
satisfied silently: the identity frame and the per-level softmax are not
housekeeping that precedes the LH-projection, they are settings whose cost
is measured before any projection is applied to them.

\medskip

This completes the apparatus. The chapter has fixed a formal setting and a
priority ordering (\S\ref{sec:methodology_problem}), derived three mechanisms that
realise that ordering on the forward map, on the assembled update and at the
branch points of a shared trunk (\S\ref{sec:methodology_hcc}--%
\S\ref{sec:methodology_lhproj}), specified the readout, decoder and
checkpoint-selection rules under which any of them is scored
(\S\ref{sec:methodology_inference}), and established that the three mechanisms
form a four-way alternative rather than a factorial
(\S\ref{sec:methodology_axis}).

Every one of those is something this thesis defines. What none of them supplies is
anything the thesis inherits: the five model families the mechanisms are applied
to, the data they are compared on, the quantities they are compared with, and the
rules under which a difference counts as evidence. \Cref{ch:experiments} supplies
those in that order, and then states the design that turns this apparatus into
answerable questions.

